\chapter{Cosmology}

\section*{Theory problems}

\probhead{13.1}{}{2008}{Bandung, Indonesia}{TH L3}{100}
% source id: 2008_TH_L3   tag: CMB horizon distance, scale factor, inflation
Measurement of the cosmic microwave background radiation (CMB) shows that its
temperature is practically the same at every point in the sky to a very high
degree of accuracy. Let us assume that light emitted at the moment of
recombination ($T_r \approx 3000$\,K, $t_r \approx 300000$ years) is only
reaching us now ($T_o \approx 3$\,K, $t_o \approx 1.5 \times 10^{10}$ years).
Scale factor $S$ is defined as such $S_0 = S(t = t_o) = 1$ and
$S_t = S(t < t_o) < 1$. Note that the radiation dominated period was between
the time when the inflation stopped ($t = 10^{-32}$ seconds) and the time when
the recombination took place, while the matter dominated period started at the
recombination time. During the radiation dominated period $S$ is proportional
to $t^{1/2}$, while during the matter dominated period $S$ is proportional to
$t^{2/3}$.

\begin{description}
\item[a.] Estimate the horizon distances when recombination took place. Assume
  that temperature $T$ is proportional to $1/S$, where $S$ is a scale factor
  of the size of the Universe.
\item[b.] Note: Horizon distance in degrees is defined as maximum separation
  between the two points in CMBR imprint such that the points could ``see''
  each other at the time when the CMBR was emitted.
\item[c.] Consider two points in CMBR imprint which are currently observed at
  a separation angle $\alpha = 5^\circ$. Could the two points communicate with
  each other using photon? (Answer with ``YES'' or ``NO'' and give the reason
  mathematically)
\item[d.] Estimate the size of our Universe at the end of inflation period.
\end{description}

\probhead{13.2}{}{2008}{Bandung, Indonesia}{TH S12}{20}
% source id: 2008_TH_S12   tag: SN Ia distance and Hubble time
Consider a type Ia supernova, in a distant galaxy, which has a luminosity of
$5.8 \times 10^{9}L_\odot$ at maximum light. Suppose you observe this
supernova using your telescope and find that its brightness is
$1.6 \times 10^{-7}$ times the brightness of Vega. The redshift of its host
galaxy is known to be $z = 0.03$. Calculate the distance of this galaxy (in
pc) using the data of the supernova and also the Hubble time.

\probhead{13.3}{}{2008}{Bandung, Indonesia}{TH S5}{20}
% source id: 2008_TH_S5   tag: dark matter/dark energy ratio at recombination
The average temperature of the Cosmic Microwave Background (CMB) is currently
$T = 2.73$\,K, and it yields the origin of CMB to be at redshift
$z_{CMB} = 1100$. The current densities of the Dark Energy, Dark Matter, and
Normal Matter components of the Universe as a whole are
$\rho_{DE} = 7.1 \times 10^{-30}\,\mathrm{g/cm^3}$,
$\rho_{DM} = 2.4 \times 10^{-30}\,\mathrm{g/cm^3}$, and
$\rho_{NM} = 0.5 \times 10^{-30}\,\mathrm{g/cm^3}$, respectively. What is the
ratio between the density of Dark Matter to the density of Dark Energy at the
time CMB was emitted, if we assume that the dark energy is vacuum energy?

\probhead{13.4}{}{2009}{Tehran, Iran}{TH P11}{10}
% source id: 2009_TH_P11   tag: neutrino number density as dark matter
Assume the mass of neutrinos is $m_\nu = 10^{-5}m_e$. Calculate the number
density of neutrinos ($n_\nu$) needed to compensate the dark matter of the
universe. Assume the universe is flat and 25\,\% of its mass is dark matter.

Hint: Take the classical total energy equal to zero.

\probhead{13.5}{}{2010}{Beijing, China}{TH S3}{10}
% source id: 2010_TH_S3   tag: QSO distance from redshift, Hubble law
The observed redshift of a QSO is $z = 0.20$, estimate its distance. The
Hubble constant is 72\,km\,s$^{-1}$\,Mpc$^{-1}$.

\probhead{13.6}{}{2011}{Katowice, Poland}{TH S11}{10}
% source id: 2011_TH_S11   tag: star-formation redshift matter-dominated
Based on the spectrum of a galaxy with redshift $z = 6.03$ it was determined
that the age of the stars in the galaxy is from 560 to 600 million years. At
what $z$ did the epoch of star formation occur in this galaxy? Assume that the
current age of the Universe is $t_0 = 13.7 \times 10^{9}$ years and that the
rate of expansion of the Universe is given by a flat cosmological model with
cosmological constant $\Lambda = 0$. (In such a model the scale factor
$R \propto t^{2/3}$, where $t$ is the time since the Big Bang.)

\probhead{13.7}{}{2011}{Katowice, Poland}{TH S15}{10}
% source id: 2011_TH_S15   tag: CMB temperature scaling with redshift
Given that the cosmic background radiation has the spectrum of a black body
throughout the evolution of the Universe, determine how its temperature
changes with redshift $z$. In particular, give the temperature of the
background radiation at the epoch $z \approx 10$ (that of the farthest
currently observed objects). The current temperature of the cosmic background
radiation is 2.73\,K.

\probhead{13.8}{}{2012}{Rio de Janeiro, Brazil}{TH S10}{}
% source id: 2012_TH_S10   tag: matter-dominated universe cooling timescale
Assume that the universe currently is well described by a density parameter
$\Omega_0 = 1$, there is no dark energy, and the current temperature of the
universe is 2.73\,K. Knowing that the temperature of the universe is inversely
proportional to its radius (the scale factor), compute how long, starting from
the present time, it will take to the Universe to cool down by 0.1\,K.

\probhead{13.9}{}{2013}{Volos, Greece}{TH L1}{50}
% source id: 2013_TH_L1   tag: critical density, galaxy escape, Hubble flow
In a homogeneous and isotropic universe, the matter (baryonic matter + dark
matter) density parameter
$\Omega_m = \dfrac{\rho_m}{\rho_c} = 32\,\%$, where $\rho_m$ is the matter
density and $\rho_c$ is the critical density of the Universe.

\begin{description}
\item[(1)] Calculate the average matter density in our local neighbourhood.
\item[(2)] Calculate the escape velocity of a galaxy 100\,Mpc away from us.
  Assume that the recession velocity of galaxies in Hubble's law equals the
  corresponding escape velocity at that distance, for the critical density of
  the Universe that we observe.
\item[(3)] The particular galaxy is orbiting around the centre of our cluster
  of galaxies on a circular orbit. What is the angular velocity of this galaxy
  on the sky?
\item[(4)] Will we ever discriminate two such galaxies that are initially at
  the same line of sight, if they are both moving on circular orbits but at
  different radii (answer ``Yes'' or ``No'')? [Assume that the Earth is
  located at the centre of our local cluster.]
\end{description}

\probhead{13.10}{}{2013}{Volos, Greece}{TH L3}{50}
% source id: 2013_TH_L3   tag: Coma cluster virial mass, dark matter M/L
\begin{description}
\item[(1)] Using the \emph{virial theorem} for an isolated, spherical system,
  i.e, that $-2\langle K\rangle = \langle U\rangle$, where ``$K$'' is the
  average kinetic energy and ``$U$'' is the average potential energy of the
  system, determine an expression for the total mass of a cluster of galaxies
  if we know the radial velocity dispersion, $\sigma$, of the cluster's galaxy
  members and the cluster's radius, $R$. Assume that the cluster is isolated,
  spherical, has a homogeneous density and that it consists of galaxies of
  equal mass.
\item[(2)] Find the \emph{virial mass}, i.e.\ the mass calculated from the
  \emph{virial theorem}, of the Coma cluster, which lies at a distance of
  90\,Mpc from us, if you know that the radial velocity dispersion of its
  member galaxies is $\sigma_{v_r} = 1000\,\mathrm{km/s}$ and that its angular
  diameter (on the sky) is about $4^\circ$.
\item[(3)] From observations, the total luminosity of the galaxies comprising
  the cluster is approximately $L = 5 \times 10^{12}\,L_\odot$. If the mass to
  luminosity ratio, $M/L$, of the cluster is ${\sim}1$ (assume that all the
  mass of the cluster is visible mass), this should correspond to a total mass
  $M \sim 5 \times 10^{12}\,\mathrm{M}_\odot$ for the mass of the cluster.
  Give the ratio of the luminous mass to the total mass of the cluster you
  derived in question (2).
\end{description}

\probhead{13.11}{}{2013}{Volos, Greece}{TH S15}{10}
% source id: 2013_TH_S15   tag: recession velocities classical vs relativistic
The Doppler shift of three remote galaxies has been measured with the help of
spectral observations:

\begin{center}
\begin{tabular}{|c|c|}
\hline
Galaxy & Redshift, $z$ \\
\hline
3C 279 & 0.536 \\
\hline
3C 245 & 1.029 \\
\hline
4C41.17 & 3.8 \\
\hline
\end{tabular}
\end{center}

\begin{description}
\item[(a)] Calculate their apparent recession velocity (1) using the classical
  approach, (2) using the approximate formula $v = c\,\ln(1+z)$, that is often
  used by cosmologists and (3) using the special relativistic approach.
\item[(b)] For all three formulae, at what percentage of the speed of light do
  they appear to recede?
\item[(c)] Which of (1) classical, (2) special relativity (3) approximate
  cosmological.
\end{description}

\probhead{13.12}{}{2013}{Volos, Greece}{TH S5}{10}
% source id: 2013_TH_S5   tag: Hubble constant and Hubble time from Halpha
The optical spectrum of a galaxy, whose distance had been measured to be
41.67\,Mpc, showed the Balmer H$\alpha$ line ($\lambda_o = 656.3$\,nm)
redshifted to $\lambda = 662.9$\,nm.
\begin{description}
\item[(a)] Use this distance to calculate a value of the Hubble constant,
  $H_o$.
\item[(b)] Using your results, estimate the Hubble time of the Universe.
\end{description}

\probhead{13.13}{}{2015}{Magelang, Indonesia}{TH S14}{}
% source id: 2015_TH_S14   tag: Olbers-type mean free path of sight line
Suppose we live in a static and infinitely large universe where the average
density of stars is $n = 10^{9}\,\mathrm{Mpc}^{-3}$ and the average stellar
radius is equal to the solar radius. Assume that standard Euclidean geometry
holds true in this universe. How far, on average, could you see in any
direction before your line of sight strikes a star? Please write your answer
in Mpc.

\probhead{13.14}{Early Universe}{2016}{Bhubaneswar, India}{TH T3}{10}
% source id: 2016_TH_T3   tag: matter-radiation equality redshift
Cosmological models indicate that radiation energy density, $\rho_r$, in the
Universe is proportional to $(1+z)^4$, and the matter energy density,
$\rho_m$, is proportional to $(1+z)^3$, where $z$ is the redshift. The
dimensionless density parameter, $\Omega$, is given as $\Omega = \rho/\rho_c$,
where $\rho_c$ is the critical energy density of the Universe. In the present
Universe, the density parameters corresponding to radiation and matter, are
$\Omega_{r0} = 10^{-4}$ and $\Omega_{m0} = 0.3$, respectively.

\begin{description}
\item[(T3.1)] Calculate the redshift, $z_e$, at which radiation and matter
  energy densities were equal. \hfill[3]
\item[(T3.2)] Assuming that the radiation from the early Universe has a
  blackbody spectrum with a temperature of 2.732\,K, estimate the
  temperature, $T_e$, of the radiation at redshift $z_e$. \hfill[4]
\item[(T3.3)] Estimate the typical photon energy, $E_\nu$ (in eV), of the
  radiation as emitted at redshift $z_e$. \hfill[3]
\end{description}

\probhead{13.15}{Super Luminal Galaxies}{2018}{Beijing, China}{TH T1}{10}
% source id: 2018_TH_T1   tag: superluminal recession true/false
Read the statements given below and state if they are true or false:

\begin{description}
\item[(a)] For some galaxies the apparent recession speed exceeds the speed of
  light.
\item[(b)] The velocity--distance relation as given by Hubble cannot allow
  recession velocities to exceed the speed of light.
\item[(c)] Hubble--Lema\^itre's law (formerly known as Hubble's Law) does not
  violate special relativity.
\item[(d)] If some galaxies would have an apparent recession speed exceeding
  the speed of light, then the photons from those galaxies can never reach us.
\item[(e)] As the expansion of the Universe is accelerating, photons emitted
  right now from galaxies which have apparent recession speed equal to the
  speed of light will never reach us.
\end{description}

\probhead{13.16}{Thermal History of the Universe}{2018}{Beijing, China}{TH T11}{75}
% source id: 2018_TH_T11   tag: Friedmann equation, CMB, equality
Based on Einstein's general relativity, Russian physicist Alexander Friedmann
derived the Friedmann Equation by which the dynamics of a homogeneous and
isotropic universe can be well described. The Friedmann Equation is usually
written as follows:
\[
\left(\frac{\dot{a}}{a}\right)^2
  = \frac{8\pi G}{3}(\rho_m + \rho_r) + \frac{\Lambda c^2}{3}
    - \frac{kc^2}{a^2}.
\]
We define the Hubble parameter as $H = \dot{a}/a$, where $a$ is the scale
factor and $\dot{a}$ is the rate of change of scale factor with time. Thus,
the Hubble parameter is a function of cosmic time. In the Friedmann Equation,
$\rho_m$ is the density of matter, including dark matter and baryons,
$\rho_r$ is the density of radiation, $\Lambda$ is the cosmological constant,
and $k$ is the curvature of space. Subscript 0 indicates the value of a
physical quantity at present day, e.g.\ $H_0$ is the present value of the
Hubble parameter. Also, to avoid confusion with the reduced Hubble parameter,
we use the reduced Planck constant $\hbar = h/(2\pi)$ instead of the Planck
constant $h$.

\begin{description}
\item[(a)] (5 points) What are the dimensions of the Hubble parameter? One can
  define a characteristic timescale for the expansion of the Universe (i.e.\
  Hubble time $t_H$) using the Hubble parameter. Calculate the present-day
  Hubble time $t_{H0}$.
\item[(b)] (5 points) Let us define the critical density $\rho_c$ as the
  matter density required to explain the expansion of a flat universe without
  any radiation or dark energy. Find an expression for the critical density,
  in terms of $H$ and $G$. Calculate the present critical density
  $\rho_{c0}$.
\item[(c)] (6 points) It is convenient to define all density parameters in a
  dimensionless manner like $\Omega_i = \rho_i/\rho_c$, i.e.\ the ratio of
  density to critical density. The Friedmann Equation can be rewritten using
  these dimensionless density parameters simply as
  \[ \Omega_m + \Omega_r + \Omega_\Lambda + \Omega_k = 1. \]
  Use this information to find expressions for $\Omega_\Lambda$ and
  $\Omega_k$, in terms of $H$, $c$, $\Lambda$, $k$ and $a$.
\item[(d)] (7 points) Another equation which is valid for matter, radiation
  and dark energy is often called the Fluid Equation:
  \[ \dot{\rho} + 3\frac{\dot{a}}{a}\left(\rho + \frac{p}{c^2}\right) = 0, \]
  where $p$ is the pressure of some component, $\rho$ is the density and
  $\dot{\rho}$ is the rate of change of density over time. Radiation contains
  photons and massless neutrinos, and they both travel at the speed of light.
  The pressure exerted by these particles is $1/3$ of their energy density.
  Show that the density of radiation $\rho_r \propto (1+z)^4$, where $z$ is
  cosmological redshift. You may note that if
  $\dot{\rho}/\rho = n\,\dot{a}/a$, then $\rho \propto a^n$.
\item[(e)] (4 points) We know that the value of the cosmological constant
  $\Lambda$ doesn't evolve. Its equation of state has the form
  $p = w\rho_\Lambda c^2$, where $w$ is an integer. Find the value of $w$.
\item[(f)] (13 points) Planck time defines a characteristic timescale before
  which our present physical laws are no longer valid, and where quantum
  gravity is needed. The expression for Planck time can be written in terms of
  $\hbar$, $G$ and $c$, and the non-dimensional coefficient of this expression
  in SI units is of the order of unity. Using dimensional analysis, find an
  expression for the Planck time and estimate its value.
\item[(g)] (7 points) The Planck length defines the length scale associated
  with the Planck time and is given by $l_P = ct_P$. The minimal mass of a
  black hole, also called the Planck mass, is defined as the mass of a black
  hole whose Schwarzschild radius is two times the Planck length. Derive the
  Planck mass $M_P$ and calculate $M_P c^2$ in GeV. This mass is considered to
  be an upper threshold for elementary particles, beyond which they will
  collapse to a black hole.
\item[(h)] (4 points) At the very beginning (soon after the Planck time), all
  the particles were in thermal equilibrium in a primordial soup. As
  temperature decreased, different particles then decoupled from the
  primordial soup one by one and could travel freely in the Universe. Photons
  decoupled at ${\sim}300\,000$ years after the Big Bang. These photons
  emitted at that time are what constitutes the cosmic microwave background
  (CMB), which follows the Stefan--Boltzmann law for blackbody radiation,
  \[ \varepsilon_r = \frac{\pi^2}{15\hbar^3 c^3}(k_B T)^4. \]
  Show that the temperature of the CMB follows $T/(1+z) = \mathrm{constant}$.
\item[(i)] (16 points) With the expansion of the Universe, radiation density
  dropped more quickly than matter density, and at some epoch the matter
  density was equal to the radiation density. Radiation contains both photons
  and neutrinos. Apart from photons, neutrinos additionally contribute to the
  radiation energy density by 68\% (i.e.\
  $\Omega_{r0} = 1.68\,\Omega_{\gamma 0}$, where $\gamma$ indicates photons).
  Estimate the redshift of matter--radiation equality $z_{eq}$ in terms of
  $\Omega_{m0}$ and the reduced Hubble parameter
  $h = H_0/(100\,\mathrm{km\,s^{-1}\,Mpc^{-1}})$. You may use the current
  temperature of the CMB: $T_0 = 2.73$\,K.
\item[(j)] (8 points) The neutrinos decoupled from the primordial soup when
  the temperature of the universe was around 1\,MeV. At this time, the
  radiation density in the universe was much more than all other components.
  Estimate the time ($t = \frac{1}{2H}$) when neutrinos decoupled, and express
  it in seconds since the big bang.
\end{description}

\probhead{13.17}{Distance to the Coma galaxy cluster}{2019}{Keszthely, Hungary}{TH 13}{40}
% source id: 2019_TH_13   tag: Hubble distance, virial mass, M/L
The Coma galaxy cluster (Abell 1656) has an angular diameter on the sky of
about 100 arcminutes, and contains more than 1000 individual galaxies, most of
which are dwarf and giant ellipticals orbiting the common center of mass of
the cluster in approximately circular orbits. The table below lists the
measured radial velocities of a few individual cluster member galaxies.

\begin{center}
\begin{tabular}{cc|cc|cc|cc}
No. & $v_r$ (km/s) & No. & $v_r$ (km/s) & No. & $v_r$ (km/s) &
No. & $v_r$ (km/s) \\
\hline
1 & 6001 &  6 & 7116 & 11 & 7156 & 16 & 7111 \\
2 & 7666 &  7 & 7004 & 12 & 7522 & 17 & 8292 \\
3 & 6624 &  8 & 4476 & 13 & 7948 & 18 & 5358 \\
4 & 5952 &  9 & 6954 & 14 & 4951 & 19 & 4957 \\
5 & 5596 & 10 & 8953 & 15 & 7797 & 20 & 7183 \\
\end{tabular}
\end{center}

\begin{description}
\item[a)] Derive the distance of the cluster from the mean radial velocity of
  the galaxies listed in the table. \hfill(8\,p)
\item[b)] Estimate the physical diameter of the cluster (in Mpc).
  \hfill(4\,p)
\item[c)] The virial theorem states that if the galaxy cluster is in dynamic
  equilibrium, then the mean kinetic energy, $\langle K\rangle$, and the mean
  gravitational potential energy, $\langle U\rangle$, are related by
  \[ -2\langle K\rangle = \langle U\rangle \]
  assuming the Coma cluster is spherical. For simplicity, assume that each
  galaxy has approximately the same mass, $m$. Use the virial theorem to prove
  that, in this case, the cluster mass $M$ (also called the virial mass) can
  be expressed as
  \[ M = \frac{5R}{G}\sigma_r^2 \]
  where $\sigma_r^2$ is the velocity dispersion of the cluster. \hfill(10\,p)

  The formula for standard deviation:
  \[ \sigma = \sqrt{\frac{\sum_{i=1}^{n}(x_i - \bar{x})^2}{n-1}} \]
\item[d)] Using the data in the table, estimate the virial mass of the Coma
  cluster in solar masses. \hfill(12\,p)
\item[e)] The total luminosity of the Coma cluster (in solar luminosity,
  $L_\odot$) is $L \approx 5 \times 10^{12}\,L_\odot$. Calculate the
  mass-luminosity ratio of the cluster in solar mass per solar luminosity
  units. \hfill(2\,p)
\item[f)] Which of the following statements is true (more than one answer is
  possible)? Mark the letter of your answer with $\times$ on the answer sheet.
  \hfill(4\,p)
  \begin{description}
  \item[(A)] The mass-luminosity ratio of the Coma cluster is much higher than
    that of a typical spiral galaxy, like the Milky Way.
  \item[(B)] The mass-luminosity ratio of the Coma cluster is similar to that
    of a typical spiral galaxy.
  \item[(C)] The mass-luminosity ratio of the Coma cluster is much less than
    that of a typical spiral galaxy.
  \item[(D)] The Coma cluster contains much more dark matter than a typical
    spiral galaxy.
  \item[(E)] The Coma cluster contains much less dark matter than a typical
    spiral galaxy.
  \end{description}
\end{description}

\probhead{13.18}{Mirror}{2020}{GeCAA (Estonia, remote)}{TH 3}{10}
% source id: 2020_TH_3   tag: H0 from laser round-trip in expanding space
A bored cosmologist comes up with a thought experiment to determine the Hubble
constant ($H_0$) for his model of a Steady-State-Universe. In this experiment,
a large, fully reflecting flat mirror -- carrying several gyroscopes that
would maintain its spatial orientation in the same plane -- would be placed at
a distance $D$ from the Solar System in a region without gravitational
influences. From the Earth, a laser beam would be directed towards that region
for a long period of time. After a time $T$, the radiation would return and be
detected, allowing the determination of the fixed constant $H_0$.

\begin{description}
\item[(a)] (7 points) Find an expression for $H_0$ as a function of $D$, $c$
  (speed of light) and $T$. Consider that the separation $S$ between the Solar
  System and the mirror increases only due to the expansion of the universe
  according to the law $S = s\,e^{H_0 t}$, where $s$ is the initial
  separation. You may use $e^x \approx 1 + x$ for $x \ll 1$, if necessary.
\item[(b)] (3 points) Imagine that such a mirror is located in the vicinity of
  the star Vega (which also features on the logo of the 1st GeCAA). Vega was
  the first star outside the Solar System to be photographed and one of the
  first stars whose parallax ($p = 0.125''$) was accurately measured in 1840
  by G.~W.~von Struve. Estimate the total duration of this $H_0$ measurement
  experiment.
\end{description}

\probhead{13.19}{Cosmic String}{2021}{Bogot\'a (remote)}{TH Q15}{55}
% source id: 2021_TH_Q15   tag: cosmic-string gravity, photon gas, lensing
\textbf{Introduction}

According to our current understanding, just after the Big Bang, when the
Universe was extremely hot, the electromagnetic force, the strong nuclear
force as well as the weak nuclear force were unified as one Grand Unified
(GUT) force.

When the Universe cooled down to $T_{GUT} = 10^{29}$\,K, the strong nuclear
force decoupled from the electroweak force. Later, when the temperature
reduced to $T_{EW} = 10^{15}$\,K, the weak force decoupled from the
electromagnetic force. These transitions happened in a rapid succession within
a small fraction of a second after the Big Bang. It is thought that these
phase transitions produced a variety of peculiar objects, called vacuum
defects, which may still be observed today.

This question will discuss properties of one such possible type of defect
called cosmic strings and their observational effects.

\emph{Note 1.} Unless otherwise stated use the laws of Newtonian Mechanics.

\emph{Note 2.} You will use the following constants:
\begin{itemize}
\item Stefan--Boltzmann constant
  \[ \sigma = \frac{2\pi^5 k_B^4}{15\,h^3 c^2}
            = \frac{\pi^2 k_B^4}{60\,\hbar^3 c^2} \]
\item The reduced Planck constant
  \[ \hbar = \frac{h}{2\pi} \]
\item Universal radiation constant
  \[ a = \frac{4\sigma}{c}
       = 7.5657 \times 10^{-16}\,\mathrm{J\,m^{-3}\,K^{-4}} \]
\item Planck temperature
  \[ T_{pl} = \sqrt{\frac{\hbar c^5}{G k_B^2}}
            = 1.416784 \times 10^{32}\,\mathrm{K} \]
\end{itemize}

\emph{Note 3.} Recall that the gravitational field $\vec{g}$ satisfies the
Gauss theorem:
\[ \oint \vec{g}\cdot d\vec{A} = -4\pi G M_{in}, \]
where $M_{in}$ is the mass enclosed by the surface $A$.

\medskip
\textbf{Part A: Gravitational Field of a Cosmic String (22 points).}

As a first approximation, let us consider a cosmic string as an infinitely
long cylinder of radius $r_0$ and mass per unit length $\mu$.

\begin{description}
\item[A.1] Write an expression in terms of the constants $G$, $\mu$ and $r_0$
  for the gravitational field produced by the string, $\vec{g}(r)$. Consider
  the cases $r_0 < r$ and $r_0 > r$ independently. \hfill[6.0\,pt]
\item[A.2] Write an expression in terms of the constants $G$, $\mu$ and $r_0$
  for $g_0 \equiv |\vec{g}(r_0)|$. \hfill[1.0\,pt]
\item[A.3] Let $g$ be defined as $\vec{g}(r)\cdot\hat{r}$. Draw a rough sketch
  of $g$ vs.\ $r$ in the figure given in the answer sheet. \hfill[3.0\,pt]
\item[A.4] It is possible to define a stable orbit around a cosmic string. For
  circular orbits of radius $R > r_0$ and period $\tau$, the following
  relation is attained
  \[ R = A\tau^\alpha, \]
  where $A$ and $\alpha$ are constants. Find $A$ and $\alpha$ in terms of $G$
  and $\mu$. \hfill[4.0\,pt]
\end{description}

The following three questions refer to a classical Newtonian particle moving
with speed $v$ when at a distance $r > r_0$ from the string. You will need to
use the result below:
\[ \int_{x_0}^{x}\frac{dx}{x} = \ln\left(\frac{x}{x_0}\right) \]

\begin{description}
\item[A.5] Show that the gravitational potential energy of the particle is
  \[ U = Gm\mu\ln\left(\frac{r}{b}\right), \]
  where $b$ is any fixed distance. \hfill[3.0\,pt]
\item[A.6] What is the maximum distance, $R_{\max}$, from the string, that the
  particle can reach? \hfill[4.0\,pt]
\item[A.7] Is it possible for the particle to escape the gravitational field?
  Write YES/NO in the answer sheet. \hfill[1.0\,pt]
\end{description}

\medskip
\textbf{Part B: Cosmic string as a photon gas (17 points).}

Consider now a cosmic string as a photon gas inside a very long cylinder of
radius $r_0$ with adiabatic walls, and in thermal equilibrium at temperature
$T$.

\begin{description}
\item[B.1] What is the energy density $\rho$ of the string in terms of $T$,
  $\hbar$, $k_B$ and $c$? \hfill[2.0\,pt]
\item[B.2] The radius $r_0$ is related to the temperature $T$ via
  \[ r_0 = \frac{\hbar^{n_1} c^{n_2}}{k_B T}, \]
  where $\hbar$ is the reduced Planck constant, $c$ is the speed of light in
  vacuum, $k_B$ is the Boltzmann constant, and $n_1$ and $n_2$ are integer
  numbers. Determine $n_1$ and $n_2$. \hfill[4.0\,pt]
\item[B.3] What is the mass per unit length, $\mu$, of the string in terms of
  $\rho$ and $r_0$? \hfill[2.0\,pt]
\item[B.4] Express the inequality for the weak field condition, defined as
  \[ \frac{2G\mu}{c^2} \ll 1, \]
  only in terms of $T$ and $T_{pl}$. \hfill[5.0\,pt]
\item[B.5] Calculate $\frac{2G\mu}{c^2}$ for
  \begin{itemize}
  \item $T = T_{EW}$
  \item $T = T_{GUT}$
  \end{itemize}
  \hfill[3.0\,pt]
\item[B.6] Does the weak field condition hold for $T_{EW}$? Answer YES or NOT.
  Does the weak field condition hold for $T_{GUT}$? Answer YES or NOT.
  \hfill[1.0\,pt]
\end{description}

\medskip
\textbf{Part C: Gravitational Lensing from cosmic Strings (16 points).}

So far, in parts A and B, we have neglected the internal pressure of the
photon gas inside the string. If we include it in our analysis, we need to
consider the General Theory of Relativity.

After solving the Einstein field equations, one finds that the spacetime
around a cosmic string is conical, as if a narrow wedge were removed from a
flat sheet and the edges connected, as shown below.

\ioaafig[1.0]{2021_TH_Q15-f4.pdf}

{\centering\small
http://www.ctc.cam.ac.uk/outreach/origins/cosmic\_structures\_five.php\par}

A remarkable result of this model is light deflection by a cosmic string,
which leads to the possibility of detection through gravitational lensing.

The angle of deflection (in radians) of a light ray coming from a distant
quasar ($O$ in the figure below), as the light passes close to a cosmic string
($S$ in the figure below) and eventually reaching an observer on the Earth
($E$ in the figure below), is
\[ \delta\phi = \frac{4\pi G\mu}{c^2} \]
and is independent of the parameter, $p$, as shown in the figure below.

In the figure $E$ and $O$ are in a plane perpendicular to the string. The
distance between the observer and the string is $D_{ES}$ and the distance
between the observer and the source is $D_{OE}$.

\ioaafig[1.0]{2021_TH_Q15-f5.pdf}
% source O offset perpendicular by distance p above the line at the far end
% (source paper Q15-5, below the deflection-angle equation)

\begin{description}
\item[C.1] Although the angle of deflection does not depend on parameter $p$,
  an Earth-based observer will be able to see more than one image only if the
  value of $p$ is within a certain range. Find a condition on the value of the
  parameter $p$ in terms of $D_{ES}$, $D_{OE}$, and temperature $T$, for an
  Earth-based observer to see more than one image of the object $O$.
  \hfill[6.0\,pt]
\item[C.2] In case the observer sees more than one image, what is the angular
  separation between each pair? Find an expression in terms of $D_{ES}$,
  $D_{OE}$ and $\delta\phi$. \hfill[6.0\,pt]
\item[C.3] If $D_{OE} = 2D_{ES}$, determine the minimum size of an optical
  telescope needed to resolve this lensing event produced by a GUT string.
  \hfill[4.0\,pt]
\end{description}

\probhead{13.20}{Journey Between Galaxies}{2022}{Kutaisi, Georgia}{TH 4}{10}
% source id: 2022_TH_4   tag: travel through expanding space, scale factor
We start a long journey to a planet located at a distance of
$d_0 = 10$\,Mpc at the start of our journey. During our travel, the universe
keeps expanding according to Hubble--Lema\^itre's law. (Assume that Hubble's
constant ($H$) does not change.)

\begin{description}
\item[(a)] Write an expression for the distance between Earth and the planet
  at a time $t$ from the start of our journey. \hfill(2\,pt)
\item[(b)] What minimum constant speed $v_0$ should the rocket have to reach
  the destination? Assuming that the speed of our rocket is
  $v_0 = 1000$\,km/s, can we reach the planet? Write YES or NO. If yes, how
  long will our journey take? \hfill(8\,pt)
\end{description}

\emph{Hint:} A simplified way to account for the expansion of our universe is
to introduce a ``scale factor'' $a(t)$, which relates the distance between two
objects $l(t)$ at a time $t$ to a distance $l_0$ at a time $t = 0$ as
$l(t) = a(t)\,l_0$.

\probhead{13.21}{Dark Energy}{2023}{Katowice, Poland}{TH Q5}{12}
% source id: 2023_TH_Q5   tag: critical density, matter-DE equality redshift
Observations indicate that the expansion of the Universe is accelerating.
Fluctuations of the cosmic microwave background favour a flat (Euclidean)
geometry, in which the total mass density (i.e.\ the sum of density of matter
and equivalent mass density of all forms of energy) should be equal to the
so-called critical density:
\[ \rho_{\mathrm{cr}} = \frac{3H_0^2}{8\pi G}, \]
where $H_0$ is the present value of the Hubble constant. However, the total
density of matter (luminous and dark) is estimated at
\[ \rho_{\mathrm{m},0} \approx 2.8 \times 10^{-27}\,\mathrm{kg\,m^{-3}}. \]
To resolve this discrepancy, the standard cosmological model assumes that the
Universe is filled with a mysterious `dark energy' of constant energy density
$\varepsilon_\Lambda$.

Determine the value of $\varepsilon_\Lambda$ and calculate for which redshift
in the past the energy density equivalent to matter was equal to the density
of dark energy. Neglect the contribution of electromagnetic radiation.

\probhead{13.22}{CMB}{2024}{Rio de Janeiro, Brazil}{TH T5}{10}
% source id: 2024_TH_T5   tag: CMB redshift, matter-dominated age
The Cosmic Microwave Background (CMB) is radiation emitted during the early
Universe. It is reasonably homogeneous and isotropic, and well described by a
black-body radiation spectrum. Its emission spectrum today has a peak at a
temperature of approximately $T_{\mathrm{today}} \sim 3$\,K, as measured by
the COBE satellite FIRAS instrument.

\begin{description}
\item[(a)] (3 points) What is the redshift ($z$) at which the CMB spectrum has
  a peak at the infrared wavelength of
  $\lambda_{\mathrm{IR}} \sim 0.1$\,mm?
\item[(b)] (7 points) By assuming a spatially flat, matter-dominated Universe,
  what is the age of the Universe corresponding to the redshift of the
  previous part?
\end{description}

\probhead{13.23}{Big Bang Nucleosynthesis}{2025}{Mumbai, India}{TH T10}{35}
% source id: 2025_TH_T10   tag: freeze-out, deuterium Saha, helium abundance
During the radiation dominated era in the early Universe, the scale factor of
the Universe $a \propto t^{1/2}$, where $t$ is the time since the Big Bang.
During most of this era, neutrons (n) and protons (p) remain in thermal
equilibrium with each other via weak interactions. The number density ($N$) of
free neutrons or protons is related to the temperature $T$ and their
corresponding masses $m$ such that
\[ N \propto m^{3/2}\exp\left(-\frac{mc^2}{k_B T}\right), \]
as long as time $t \leq t_{\mathrm{wk}} = 1.70$\,s, when
$k_B T \geq k_B T_{\mathrm{wk}} = 800$\,keV. After $t_{\mathrm{wk}}$, the weak
interactions can no longer maintain such equilibrium, and free neutrons decay
to protons with a half-life time of $610.4$\,s.

\begin{description}
\item[(T10.1)] Let the number density of protons be $N_{\mathrm{p}}$, and that
  of neutrons be $N_{\mathrm{n}}$. Calculate the relative abundance of
  neutrons given the ratio
  $X_{\mathrm{n,wk}} = N_{\mathrm{n}}/(N_{\mathrm{n}} + N_{\mathrm{p}})$ at
  time $t_{\mathrm{wk}}$. \hfill[4]
\item[(T10.2)] Photons maintain thermal equilibrium and retain a blackbody
  spectrum at all epochs.
  \begin{description}
  \item[(T10.2a)] Find the index $\beta$, such that $T(a) \propto a^\beta$.
    \hfill[2]
  \item[(T10.2b)] Identify which of the following graphs shows the correct
    behaviour of the spectral energy density for two temperatures $T_1$ and
    $T_2$. Tick the correct option in the Summary Answersheet. \hfill[2]
  \end{description}

  \begin{center}
  \includegraphics[width=0.48\linewidth]{2025_TH_T10-f1-panelA.pdf}\hfill
  \includegraphics[width=0.48\linewidth]{2025_TH_T10-f1-panelB.pdf}\par\smallskip
  \includegraphics[width=0.48\linewidth]{2025_TH_T10-f1-panelC.pdf}\hfill
  \includegraphics[width=0.48\linewidth]{2025_TH_T10-f1-panelD.pdf}\par\smallskip
  \includegraphics[width=0.48\linewidth]{2025_TH_T10-f1-panelE.pdf}
  \end{center}

\item[(T10.3)] After $t_{\mathrm{wk}}$, the process of formation of deuterium
  from protons and neutrons is governed by the Saha equation, given by the
  Indian physicist Prof.\ Meghnad Saha, which can be simplified to
  \[ \frac{N_{\mathrm{D}}}{N_{\mathrm{n}}}
     = 6.5\,\eta\left(\frac{k_B T}{m_{\mathrm{n}}c^2}\right)^{3/2}
       \exp\left(-\frac{(m_{\mathrm{D}} - m_{\mathrm{p}}
       - m_{\mathrm{n}})c^2}{k_B T}\right). \]
  Here, the baryon-to-photon ratio $\eta$ is $6.1 \times 10^{-10}$, and
  $N_{\mathrm{D}}$ is the number density of deuterium.
  \begin{description}
  \item[(T10.3a)] Plot the ratio $N_{\mathrm{D}}/N_{\mathrm{n}}$ on the grid
    in the Summary Answersheet, for at least 4 reasonably spaced values of
    temperature that lie in the domain $k_B T = [60, 70]$\,keV, and draw a
    smooth curve passing through these points. \hfill[5]
  \item[(T10.3b)] From the plot find $k_B T_{\mathrm{nuc}}$ (in keV) when
    $N_{\mathrm{D}} = N_{\mathrm{n}}$. \hfill[1]
  \item[(T10.3c)] Instead, now assume that all the free neutrons combine
    instantaneously with the protons at $k_B T_{\mathrm{nuc}}$ to form
    deuterium, all of which immediately gets converted to helium
    ($^4_2$He). Compute the corresponding epoch or time of nucleosynthesis,
    $t_{\mathrm{nuc}}$ (in s), for the formation of helium. \hfill[4]
  \end{description}
\item[(T10.4)] Calculate the value of $X_{\mathrm{n,nuc}}$ immediately before
  $t_{\mathrm{nuc}}$. \hfill[5]
\item[(T10.5)] The primordial helium abundance, $Y_{\mathrm{prim}}$, is
  defined to be the fraction of total mass in the Universe that is bound in
  helium at $T_{\mathrm{nuc}}$. Obtain a theoretical estimate for the value of
  $Y_{\mathrm{prim}}$. For the purpose of this calculation alone, assume
  $m_{\mathrm{p}} \approx m_{\mathrm{n}}$ and that the mass of helium,
  $m_{\mathrm{He}} \approx 4m_{\mathrm{n}}$. \hfill[3]
\item[(T10.6)] The primordial abundance of helium is very difficult to
  measure, as stars continuously convert hydrogen to helium in the Universe.
  The amount of processing by stars in a galaxy is characterised by the
  relative number density of oxygen (which is only produced by stars) to
  hydrogen, denoted as (O/H), in the galaxy. A compilation of the measurements
  of (O/H) and the helium abundance, $Y$, for different galaxies is plotted
  below. Use all of the points in this plot (which is reproduced in the
  Summary Answersheet) to answer the following.

  \ioaafig{2025_TH_T10-f2.pdf}

  \begin{description}
  \item[(T10.6a)] Estimate $Y$ for a blue compact dwarf galaxy with a value of
    $\mathrm{(O/H)} = 1.75 \times 10^{-4}$. \hfill[2]
  \item[(T10.6b)] Obtain the slope $dY/d\mathrm{(O/H)}$ of the straight line
    fit to the above data. \hfill[2]
  \item[(T10.6c)] Estimate the primordial helium abundance,
    $Y^{\mathrm{obs}}_{\mathrm{prim}}$, based on the above observations.
    \hfill[2]
  \end{description}
\item[(T10.7)] The deviation between $Y_{\mathrm{prim}}$ and
  $Y^{\mathrm{obs}}_{\mathrm{prim}}$ can be reconciled by changing the
  baryon-to-photon ratio $\eta$. When $\eta$ is decreased, as indicated by
  $\downarrow$ in the Summary Answersheet, indicate the increase ($\uparrow$)
  or decrease ($\downarrow$) in $N_{\mathrm{D}}/N_{\mathrm{n}}(T)$,
  $T_{\mathrm{nuc}}$ (when $N_{\mathrm{D}} = N_{\mathrm{n}}$),
  $t_{\mathrm{nuc}}$, $X_{\mathrm{n,nuc}}$, and $Y_{\mathrm{prim}}$ in the
  boxes provided in the Summary Answersheet. \hfill[3]
\end{description}

\section*{Data-analysis problems}

\probhead{13.24}{Virgo Cluster}{2008}{Bandung, Indonesia}{DA D1}{100}
% source id: 2008_DA_D1   tag: Virgo distance, H0, Hubble time
The Virgo cluster of galaxies is the nearest large cluster which extends over
nearly 10 degrees across the sky and contains a number of bright galaxies. It
is very interesting to find the distance to Virgo and to deduce certain
cosmological information from it. The table below provides the distance
estimates using various distance indicators (listed in the left column). The
right column lists the mean distance $d_i \pm$ the standard deviation $s_i$.

\begin{center}
\begin{tabular}{cl c}
$i$ & Distance Indicator & Virgo Distance (Mpc) \\
\hline
1 & Cepheids                       & $14.9 \pm 1.2$ \\
2 & Novae                          & $21.1 \pm 3.9$ \\
3 & Planetary Nebulae              & $15.2 \pm 1.1$ \\
4 & Globular Cluster               & $18.8 \pm 3.8$ \\
5 & Surface Brightness Fluctuation & $15.9 \pm 0.9$ \\
6 & Tully-Fisher relation          & $15.8 \pm 1.5$ \\
7 & Faber-Jackson relation         & $16.8 \pm 2.4$ \\
8 & Type Ia Supernovae             & $19.4 \pm 5.0$ \\
\end{tabular}
\end{center}

\begin{description}
\item[1.] By applying a weighted mean, compute the average distance (which can
  be taken as an estimate to the distance to Virgo)
  \[ d_{\mathrm{avg}}
     = \frac{\sum_i \dfrac{d_i}{s_i^2}}{\sum_i \dfrac{1}{s_i^2}} \]
  where the sum runs over the eight distance indicators used.
\item[2.] What is the uncertainty (rms) (in units of Mpc) in that estimate?
\item[3.] Spectra of the galaxies in Virgo indicate an average recession
  velocity of 1136\,km/sec for the cluster. Can you estimate the Hubble
  constant $H_0$ and its uncertainty (rms)?
\item[4.] What is the Hubble Time (age of the universe) using the value of the
  Hubble constant you found and the uncertainty (rms)?
\end{description}

\probhead{13.25}{The search for dark matter}{2017}{Phuket, Thailand}{DA D2}{75}
% source id: 2017_DA_D2   tag: LSB rotation curve, DM halo profile fit
A low surface brightness galaxy (LSB) is a diffuse galaxy with a surface
brightness that, when viewed from the Earth, is at least one magnitude lower
than the ambient night sky.

Some of its matter is in the form of ``baryonic'' matter such as neutral
hydrogen gas and stars. However, most of its matter is in the form of
invisible mass -- so called ``dark matter''. In this question, we will
investigate the mass of dark matter in a galaxy, the effect of dark matter on
the rotation curves of the galaxy, and the distribution of dark matter in the
galaxy.

The table below provides the data of a LSB galaxy named UGC4325. The galaxy is
assumed to be edge-on. At every distance $r$ from the centre of the galaxy, we
measure
\begin{enumerate}
\item $\lambda_{\mathrm{obs}}$, the observed wavelength of the H$\alpha$
  emission line. The Hubble expansion of the Universe has already been
  excluded from the data.
\item $V_{\mathrm{gas}}$, the contribution of the gas component to the
  rotation due to $M_{\mathrm{gas}}$, derived from HI surface densities.
\item $V_*$, the contribution of the stellar component to the rotation due to
  $M_*$, derived from R-band photometry.
\end{enumerate}
The rotational velocities of the test particle due to the gas component,
$V_{\mathrm{gas}}$, and the star component, $V_*$, are defined as the
velocities in the plane of the galaxy that would result from the corresponding
components without any external influences. These velocities are calculated
from the observed baryonic mass density distributions.

\begin{center}
\begin{tabular}{cccc}
$r$ & $\lambda_{\mathrm{obs}}$ & $V_{\mathrm{gas}}$ & $V_*$ \\
(kpc) & (nm) & (km/s) & (km/s) \\
\hline
 0.70 & 656.371 &  2.87 & 20.97 \\
 1.40 & 656.431 &  6.75 & 32.22 \\
 2.09 & 656.464 & 14.14 & 40.91 \\
 2.79 & 656.475 & 20.18 & 46.75 \\
 3.49 & 656.478 & 24.08 & 50.10 \\
 4.89 & 656.484 & 28.08 & 47.94 \\
 6.25 & 656.481 & 29.25 & 45.47 \\
 7.10 & 656.481 & 27.03 & 47.78 \\
 9.03 & 656.482 & 25.90 & 45.32 \\
12.05 & 656.482 & 21.03 & 42.30 \\
\end{tabular}
\end{center}

The mass of dark matter $M_{\mathrm{DM}}(r)$ within a volume of radius $r$ can
be defined in terms of the rotational velocity due to dark matter
$V_{\mathrm{DM}}$, the radius $r$ and the gravitational constant $G$,
\begin{equation}
M_{\mathrm{DM}}(r) = \frac{rV_{\mathrm{DM}}^2}{G}. \tag{1}
\end{equation}
To a good approximation, the observed rotational velocity $V_{\mathrm{obs}}$
can be modelled as
\begin{equation}
V_{\mathrm{obs}}^2 = V_{\mathrm{gas}}^2 + V_*^2 + V_{\mathrm{DM}}^2.
\tag{2}
\end{equation}
The observed rotational velocity $V_{\mathrm{obs}}$ depends on the mass of the
galaxy $M(r)$ within a volume of radius $r$ measured from the galaxy's centre.

The mass density $\rho_{\mathrm{DM}}(r)$ of dark matter within a volume of
radius $r$ can be modelled by a galaxy density profile,
\begin{equation}
\rho_{\mathrm{DM}}(r)
  = \frac{\rho_0}{1 + \left(\dfrac{r}{r_{\mathrm{C}}}\right)^2}
\tag{3}
\end{equation}
where $\rho_0$ and $r_{\mathrm{C}}$ are the central density and the core
radius of the galaxy, respectively. According to the density profile, the mass
of dark matter $M_{\mathrm{DM}}(r)$ within a volume of a radius $r$ can be
described by
\begin{equation}
M_{\mathrm{DM}}(r) = 4\pi\rho_0 r_{\mathrm{C}}^2
  \left[r - r_{\mathrm{C}}\arctan(r/r_{\mathrm{C}})\right].
\tag{4}
\end{equation}

\medskip
\textbf{Part 1 The mass of dark matter and rotation curves of the galaxy}

\begin{description}
\item[(D2.1)] In laboratories on Earth, H$\alpha$ has an emitted wavelength
  $\lambda_{\mathrm{emit}}$ of 656.281\,nm. Compute the observed rotational
  velocities of the galaxy $V_{\mathrm{obs}}$ and the rotational velocities
  due to the dark matter $V_{\mathrm{DM}}$ at distance $r$ in units of
  km\,s$^{-1}$.

  For the different values of $r$ given in the table, compute the dynamical
  mass $M(r)$ and the mass of dark matter $M_{\mathrm{DM}}(r)$ in solar
  masses. \hfill[21]
\item[(D2.2)] Create rotation curves of the galaxy on graph paper by plotting
  the points of $V_{\mathrm{obs}}$, $V_{\mathrm{DM}}$, $V_{\mathrm{gas}}$,
  $V_*$ versus the radius $r$ and draw smooth curves through the points (mark
  your graph as ``D2.2'').

  Order the contribution of the different components to the observed velocity
  in descending order. \hfill[16]
\end{description}

\medskip
\textbf{Part 2 Dark matter distribution}

\begin{description}
\item[(D2.3)] Take a data point at small $r$ and large $r$ to estimate
  $\rho_0$ and $r_{\mathrm{C}}$. Note that for large values of $x$,
  $\arctan(x) \approx \pi/2$ and at small $x$,
  $\arctan(x) \approx x - x^3/3$. \hfill[7]
\item[(D2.4)] By comparing Equation (4) to a linear function, the central
  density $\rho_0$ could also be found by a linear fit. Plot an appropriate
  graph so that a linear fit can be used to find another value of $\rho_0$.
  Evaluate $\rho_0$ in units of $M_\odot$\,kpc$^{-3}$. (Mark your graph as
  ``D2.4''). If you cannot find the value of $r_{\mathrm{C}}$ from the
  previous part, use $r_{\mathrm{C}} = 3.2$\,kpc as an estimate for this
  part. \hfill[19]
\item[(D2.5)] Compute logarithmic values of the dark matter density,
  $\ln[\rho_{\mathrm{DM}}(r)]$, and plot the distribution of dark matter in
  the galaxy as a function of radius $r$ on graph paper. (Mark your graph as
  ``D2.5''). \hfill[12]
\end{description}

\probhead{13.26}{Scaling Relations}{2021}{Bogot\'a (remote)}{DA DA1}{75}
% source id: 2021_DA_DA1   tag: Tully-Fisher fit, dark-to-baryonic masses
Spiral galaxies are disk-like rotating structures, whose dynamical state is
fairly grasped by the so-called rotation curves, quantifying the mean
rotational velocity of the disk at different distances from the center (see
Figure 1, curve B). An interesting feature is the flat region of the curve,
which is attributed to the presence of dark matter. Without it, rotation
velocities would drop steadily at large radii from the center, as depicted in
curve A.

\ioaafig[1.0]{2021_DA_DA1-f5.pdf}
% vs radius: curve A falling at large radii, curve B flattening (source paper
% Q1-1, "Figure 1: Rotation curves")

In disk galaxies a strong correlation has been observed between the intrinsic
luminosity of the whole galaxy and the asymptotic rotational velocity (as
given by the rotation curve for the outer edge of the galaxy, i.e.\
$R_{\max}$), a result that is known as the Tully--Fisher relation. This
relation also holds if you use the luminosity in a specific band. This is
shown on Figure 2 for a number of galaxies in a galaxy cluster. Every dot is a
galaxy, and the solid line is the best-fit linear relation between absolute
magnitude in $K$ band and $\log_{10}(V_{\max})$ for the whole sample.

\ioaafig{2021_DA_DA1-f1.pdf}

{\centering\small Figure 2: Absolute magnitude in $K$ band vs
$\log_{10}(V_{\max}\,[\mathrm{km\,s^{-1}}])$. Tully--Fisher relation for
several galaxies. Every dot represents a galaxy. The dark points are five
selected galaxies, for which we will provide some numbers in part 1.2.\par}

Another interesting trend is shown in Figure 3: disks with larger stellar
masses ($M_*$) tend to have smaller gas fractions ($M_{gas}/M_*$).

\ioaafig{2021_DA_DA1-f2.pdf}

In the following questions you will be asked to extract physical information
about the galaxies using the scaling relations just introduced. Consider the
following guidelines:
\begin{itemize}
\item Assume that $V_{\max}$ was measured at the same radius for all galaxies
  ($R_{\max}$), in the flat part of the rotation curves and well beyond the
  end of the stellar disk.
\item Use $M_{dm}$ for the dark matter mass up to $R_{\max}$ and $M_{tot}$
  for the sum of all components (gas, stars and dark matter).
\item Assume that all galaxies have identical stellar populations (the term
  stellar population refers to the type of stars that are present in a
  galaxy, and the relative amount of each type with respect to the total
  number of stars), and assume that the gaseous component does not interact
  with the stellar light.
\item The galaxy cluster is far away. Its distance is much larger than the
  cluster size.
\item In spherically-symmetric mass distributions, to infer the gravitational
  effect on a particle at distance $r$ from the center, it suffices to
  consider the total mass enclosed up to that radius $M(\leq r)$ as if it
  were placed at the very center of the distribution.
\end{itemize}

\medskip
\textbf{Part 1 (20 points).}

\begin{description}
\item[1.1] From an analysis of Figure 3, find the appropriate constants in the
  following relation: $M_{gas} = a \times M_*^b$
  \[ a = ?\qquad b = ? \]
  \hfill[5.0\,pt]
\item[1.2] In the plot of the Tully--Fisher relation there are 5 highlighted
  points. Data for these 5 galaxies is given in the following table. Use this
  dataset to find the appropriate constants for the TF relation presented
  below the table, by means of a linear fit using the method of least
  squares.

  \emph{Note:} Treat $\log_{10}(V_{\max})$ as the $x$ variable and $K$ as the
  $y$ variable in the linear fit.

  \begin{center}
  \begin{tabular}{cc}
  $V_{\max}$ [km/s] & $K$ [mag] \\
  \hline
   79.4 & $-16.8$ \\
  100.1 & $-19.2$ \\
  158.5 & $-21.3$ \\
  251.2 & $-21.4$ \\
  316.2 & $-24.0$ \\
  \end{tabular}
  \end{center}

  \[ K = c \times \log_{10}(V_{\max}) + d \]
  \[ c = ?\qquad d = ? \]
  \hfill[15.0\,pt]
\end{description}

\medskip
\textbf{Part 2 (16 points).}

For two galaxies, G1 and G2, in the cluster, the recorded \emph{apparent}
magnitudes are:
\[ k_1 = 19.2\,;\qquad k_2 = 25.2 \]
Using this information and the relations calibrated in Part 1, find the
correct exponents in the following equations:

\begin{description}
\item[2.1] $\dfrac{M_{*1}}{M_{*2}} = 10^e$;\quad $e = ?$ \hfill[6.0\,pt]
\item[2.2] $\dfrac{M_{gas1}}{M_{gas2}} = 10^f$;\quad $f = ?$ \hfill[4.0\,pt]
\item[2.3] $\dfrac{M_{tot1}}{M_{tot2}} = 10^g$;\quad $g = ?$ \hfill[6.0\,pt]
\end{description}

\medskip
\textbf{Part 3 (15 points).}

\begin{description}
\item[3.1] Fill in the missing values in the table using the fact that for
  galaxy $G_1$, the dark-to-baryonic mass ratio up to $R_{\max}$ is 6.82.
  \hfill[15.0\,pt]

  \begin{center}
  \begin{tabular}{cccccc}
  Galaxy & Apparent magnitude $k$ & $M_{gas}\,[M_\odot]$ & $M_*\,[M_\odot]$ &
  $M_{dm}\,[M_\odot]$ & $M_{tot}\,[M_\odot]$ \\
  \hline
  $G_1$ & 19.2 & & & & $4.39 \times 10^{11}$ \\
  \end{tabular}
  \end{center}
\end{description}

\medskip
\textbf{Part 4 (24 points).}

\begin{description}
\item[4.1] Consider a systematic uncertainty of $\sigma_{sys} = \pm 0.2$ in
  each apparent magnitude due to CCD calibration errors. Then $k_1$ must be
  read as $k_1 = 19.2 \pm 0.2$, i.e., the only thing we know is that $k_1$
  most likely lies in the interval $[19.0, 19.4]$. The same goes for $k_2$.

  Recalculate the exponent in the scaling relation
  $\frac{M_{*1}}{M_{*2}} = 10^e$ (found in 2.1), expressing $e$ as an
  interval estimated by considering the extreme possible variations in $k_1$
  and $k_2$.
  \[ e \in [?, ?] \]
  \hfill[4.0\,pt]
\item[4.2] Now we consider that there is always a natural spread of the data
  around any relation. For instance, for a given value of the $K$ magnitude
  the TF relation gives a single value of $\log_{10}(V_{\max})$, but it would
  be more realistic to report an interval of plausible values, derived from
  the natural spread of the data around the mean TF relation. We call this
  the statistical uncertainty, $\sigma_{stat}$.

  Estimate the statistical uncertainty if $\log_{10}(V_{\max})$ is inferred
  from $K$ using the TF relation from question 1.2. For this, consider for
  each point the difference between the value of $\log_{10}(V_{\max})$
  estimated from $K$ using your linear fit and the actual measurement of
  $\log_{10}(V_{\max})$, and take $\sigma_{stat}$ as two times the root mean
  square (RMS) of these differences (the RMS of a set of values is the square
  root of the arithmetic mean of the squares of those values).
  \[ \sigma_{stat} = ? \]
  \hfill[10.0\,pt]
\item[4.3] Recalculate the exponent in the scaling relation
  $\frac{M_{tot1}}{M_{tot2}} = 10^g$, expressing $g$ as an interval estimated
  by considering the extreme possible variations arising from both the
  systematic and statistical uncertainties:
  \[ g \in [?, ?] \]
  \hfill[10.0\,pt]
\end{description}

\probhead{13.27}{Galactic Surveys}{2022}{Kutaisi, Georgia}{DA 2}{105}
% source id: 2022_DA_2   tag: halo pair statistics, redshift-space distortions
In Cosmology, observations of the spatial distribution of galaxies are crucial
to understanding the evolution and expansion of the Universe. Numerous
telescopes conduct surveys to produce catalogues of galaxies, which are then
subjected to statistical analysis. Usually, galaxies form in groups, which
together are part of a cluster, surrounded by a very large, spherical
dark-matter halo. Most of the matter is concentrated in the centre, thus there
is a high probability of having a large-mass galaxy at the centre of the halo.

In this problem, we will assume that all halos have one central galaxy and
several groups of satellite galaxies, and use statistical analysis to draw
inferences about the spatial distribution of galaxies. In Figure 1, you can
see six different possible probability densities (A--F) for distances between
galaxies. On the X axis is the separation distance between galaxies ($a$), and
on the Y axis is the probability density ($f_p(a)$). The probability of the
distance between a pair of galaxies being in the range $a$ to $a + \Delta a$
is given by
\[ P(a) = f_p(a)\,\Delta a \]
Note that $f_p(a)$ may take values greater than 1 in some cases.

\ioaafig{2022_DA_2-f5.pdf}
% densities f_p(a) vs separation a (source paper, Data Analysis page 3 of 7,
% "Figure 1: Probability density of all possible galaxy pairs")

\medskip
\textbf{Part 1 (60\,pt)}

Table 1 gives the observational data of a single dark matter halo in the form
of Cartesian coordinates of 16 satellite groups and the coordinates of the
central galaxy, which is given in the first row. The $i^{\mathrm{th}}$ group
consists of $N_i$ galaxies. The origin of this coordinate system is the Earth.
For simplicity we assume that a group of galaxies means a collection of
galaxies with almost the same coordinates.

\begin{center}
\begin{tabular}{ccccc}
$i$ & $x$ (Mpc) & $y$ (Mpc) & $z$ (Mpc) & $N_i$ \\
\hline
C  & 8.300 & 6.200 &  1.100 & Central \\
\hline
1  & 8.401 & 6.309 &  1.394 & 291 \\
2  & 9.883 & 7.189 &  1.506 & 170 \\
3  & 8.883 & 6.413 &  2.226 & 253 \\
4  & 8.444 & 6.569 &  2.439 & 8   \\
5  & 7.782 & 6.048 & $-0.358$ & 72  \\
6  & 8.780 & 6.305 &  2.463 & 120 \\
7  & 7.990 & 5.881 &  0.532 & 302 \\
8  & 8.540 & 6.388 &  0.369 & 52  \\
9  & 7.975 & 6.030 &  1.216 & 28  \\
10 & 7.072 & 5.431 &  0.488 & 40  \\
11 & 8.037 & 5.965 &  0.615 & 72  \\
12 & 8.681 & 6.483 &  0.734 & 379 \\
13 & 7.115 & 5.259 &  1.403 & 82  \\
14 & 9.587 & 7.130 &  0.322 & 62  \\
15 & 8.193 & 6.104 &  0.915 & 305 \\
16 & 7.613 & 5.847 &  1.666 & 293 \\
\end{tabular}
\end{center}

\begin{description}
\item[(a)] The distribution of the satellite groups in the $X-Z$ plane and
  $Y-Z$ plane has been plotted for you. Plot the distribution in the $X-Y$
  plane on millimeter paper. \hfill(6\,pt)

  \ioaafig{2022_DA_2-f2.pdf}

\item[(b)] Given that this cluster of galaxies is shaped like a disk and our
  line of sight is in the plane of the disk, find the unit vector in the
  direction of the normal to the disk. \hfill(4\,pt)
\item[(c)] Let us define a new coordinate system centered on the central
  galaxy, defining our line of sight from the Earth to the central galaxy as
  the $X'$ axis and the normal to the disk as the $Z'$ axis. Find the unit
  vectors for this coordinate system. \hfill(2\,pt)
\item[(d)] Convert the given Cartesian coordinates to the new $X'Y'Z'$
  coordinates. Calculate $X'$, $Y'$, $Z'$. \hfill(24\,pt)
\item[(e)] Write the equations necessary to convert the new $X'Y'Z'$
  Cartesian coordinates into spherical coordinates centered on the central
  galaxy. The $\theta$ coordinate has been already calculated for you below.
  Calculate $R$ and $\phi$. \hfill(11\,pt)

  \begin{center}
  \begin{tabular}{cc}
  $i$ & $\theta$ \\
  \hline
  1  & 1.652 \\
  2  & 1.490 \\
  3  & 1.432 \\
  4  & 1.721 \\
  5  & 1.692 \\
  6  & 1.430 \\
  7  & 1.474 \\
  8  & 1.580 \\
  9  & 1.723 \\
  10 & 1.646 \\
  11 & 1.519 \\
  12 & 1.569 \\
  13 & 1.542 \\
  14 & 1.557 \\
  15 & 1.516 \\
  16 & 1.705 \\
  \end{tabular}
  \end{center}

  \emph{Note:} in spherical coordinates, one uses the coordinates $R$,
  $\theta$ and $\phi$, where $R$ is the radial distance, $\phi$ is the
  azimuthal angle measured anti-clockwise from the $X'$ axis (from 0 to
  $2\pi$) and $\theta$ is the polar angle (from 0 to $\pi$), i.e.\ the angle
  relative to the $Z'$ axis.
\item[(f)] Plot a histogram (on millimeter paper) showing the number of
  satellite galaxies at different radial distances from the central galaxy.
  Use 8 equally sized bins for radial distances, with 0.2\,Mpc to 2\,Mpc as
  the total range of the histogram. \hfill(9\,pt)
\item[(g)] The histogram must also be converted into a probability density. On
  the right hand side of the vertical axis of the histogram you plotted, add
  markings for the re-scaled vertical axis appropriate for the probability
  density. You will use this probability density data in part 3. \hfill(4\,pt)
\end{description}

\medskip
\textbf{Part 2 (17\,pt)}

The coordinates in the Table above are determined using distances obtained by
the redshift method. It is well-known that while measuring distance to any
galaxy with the redshift method, errors might occur because of the orbital
motion of the galaxies. Now, we will try to correct for these errors. Assume
that the central galaxy has an effective mass of $1 \times 10^{10}\,M_\odot$,
all satellite galaxies move in circular orbits around the central galaxy and
interactions between any two satellite galaxies can be ignored.

\begin{description}
\item[(a)] Use the data in Table 1 of the question sheet to obtain a histogram
  of number of galaxies against the $\phi$ azimuthal angle (take the width of
  the bins as $60^\circ$). The angle distribution is expected to be uniform,
  but it is not. Calculate the standard deviation for this data. \hfill(5\,pt)
\item[(b)] The two bins centered along the $Y'$ axis show an unusually low
  number of galaxies as the orbital velocities of satellite galaxies in these
  bins are almost parallel to our line of sight and hence there is maximum
  error in their distance estimated by the redshift method. Derive an
  expression for the correction in each azimuthal angle. \hfill(4\,pt)

  \emph{Note:} You may assume that the distance from the central galaxy ($R$)
  doesn't change significantly because of this correction.
\item[(c)] Calculate the corrected value of the azimuthal angles for each data
  row and recalculate the histogram and its standard deviation of its values.
  \hfill(8\,pt)
\end{description}

\medskip
\textbf{Part 3 (28\,pt)}

So far, we have only calculated the central galaxy -- satellite galaxy
histogram for a single halo. It is time to move on to calculating other
relations. The same observations and data analysis has been carried out for
different halos and their central galaxies. Suppose we have 500 halos, thus,
500 central galaxies ($N_C = 500$). Take that the average number of satellite
galaxies in each halo is 1000 ($N_S = 1000$).

\begin{description}
\item[a)] Calculate the following quantities:
  \begin{description}
  \item[(1)] $N_{CC}$, total number of pairs of central galaxies.
  \item[(2)] $N_{CS}^*$, total number of central galaxy -- satellite pairs
    such that the galaxies in the pair belong to the same halo.
  \item[(3)] $N_{CS}$, total number of central galaxy -- satellite pairs such
    that the galaxies in the pair belong to two separate halos.
  \item[(4)] $N_{SS}^*$, total number of satellite -- satellite galaxy pairs
    such that the galaxies in the pair belong to the same halo.
  \item[(5)] $N_{SS}$, total number of satellite -- satellite galaxy pairs
    such that the galaxies in the pair belong to two separate halos.
  \end{description}

  The table below gives the probability density values for different distances
  between central galaxies in different dark matter halos (CC). You should
  note this data is not real and doesn't represent the results of the real
  observations. Dark matter halos do not intersect. But for simplicity we will
  use this data.

  \begin{center}
  \begin{tabular}{cc}
  $R$ (Mpc) & $\rho_{CC}$ \\
  \hline
  0.3125 & 0.406 \\
  0.5375 & 0.808 \\
  0.7625 & 1.341 \\
  0.9875 & 1.134 \\
  1.2125 & 0.479 \\
  1.4375 & 0.143 \\
  1.6625 & 0.073 \\
  1.8875 & 0.060 \\
  \end{tabular}
  \end{center}

  We can calculate all probability densities using the given Center-Center
  (CC) data and the same halo Center-Satellite ($CS^*$) data that we have
  calculated in part 1. We use the following relations:
  \[ \rho_{SS}^*(R) = c_1 R\left(\rho_{CS}^*(R)\right)^2 \]
  \[ \rho_{CS}(R) = c_2 R\left(\rho_{CS}^*(R)\right)^2
     \sqrt{\rho_{CC}(R) + 1} \]
  \[ \rho_{SS}(R) = c_3 R\left(\rho_{CS}^*(R)\left(\rho_{CC}(R)
     - 5\right)\right)^2 \]
  where $c_1$, $c_2$ and $c_3$ are normalization constants and should be found
  for each probability density.
\item[b)] Calculate normalized $\rho_{SS}^*(R)$, $\rho_{CS}(R)$,
  $\rho_{SS}(R)$. \hfill(15\,pt)
\item[c)] Calculate the final probability density of all galaxy pairs, that
  combines all the distributions above with appropriate weights. Which one of
  the 6 graphs on Fig.~1 best represents this probability density?
  \hfill(8\,pt)
\end{description}

\clearpage
\section*{Solutions to Chapter 13}

\solhead{13.1}{}
% source id: 2008_TH_L3
\ioaafig{2008_TH_L3-s1.pdf}

\begin{description}
\item[a)] $X_0 = d\theta_0$, where $d = c(t_o - t_r) \approx ct_o$,
$t_o \gg t_r$.

Scale factor $S$ is defined as such $S_0 = S(t=t_o) = 1$ and
$S_t = S(t<t_o) < 1$. Maximum distance between objects at the time of
recombination
\[
X_r = S_r X_o = S_r \theta_o t_o c
\]
On the other hand $X_r = ct_r$, such that
\[
ct_r = S_r\theta_0 t_o c
\qquad\text{or}\qquad
\theta_o = \frac{t_r}{S_r t_o}
\]
\hfill(20\%)

From $T \propto \dfrac{1}{S}$ then
$\dfrac{S_o}{S_r} = \dfrac{T_r}{T_o}$ and
$\dfrac{1}{S_r} = \dfrac{T_r}{T_o}$.

So,
\[
\theta_o = \frac{t_r T_r}{t_o T_o}
= \frac{(3\times10^{5}\ yr)(3000\ K)}{(1.5\times10^{10}\ yr)(3\ K)}
= 2\times10^{-2}\ rad = 1.14^\circ
\]
\hfill(20\%)

\item[b)] NO, because $\alpha > \theta_o$
\hfill(20\%)

\item[c)]
\ioaafig{2008_TH_L3-s2.pdf}
From the time at recombination ($t_r$) to the current time ($t_0$)
$S \propto t^{2/3}$, we derive
\[
\frac{S_o}{S_r} = \left(\frac{t_o}{t_r}\right)^{2/3}
= \left(\frac{1.5\times10^{10}\ years}{3\times10^{5}\ years}\right)^{2/3}
= 1.36\times10^{3}
\]
The size of our universe at $t_r$ is
\[
\frac{1.5\times10^{10}\ \text{light years}}{1.36\times10^{3}}
= 1.11\times10^{7}\ \text{light years}
\]
\hfill(20\%)

When the inflation stopped, the Universe was radiation dominated
($t_r = 3\times10^{5}$ years to $t_i = 10^{-32}$ s), we get
\[
\frac{S_r}{S_i} = \left(\frac{t_r}{t_i}\right)^{1/2}
= \left(\frac{9.46\times10^{12}\ s}{10^{-32}\ s}\right)^{1/2}
= 3.08\times10^{22}.
\]
The size of universe at the end of the inflation period was
\[
\frac{1.1\times10^{7}\ \text{light years}}{3.08\times10^{22}}
= 3.57\times10^{-16}\ \text{light years} \approx 3\ m
\]
\hfill(20\%)
\end{description}

\solhead{13.2}{}
% source id: 2008_TH_S12
The relation between flux $F$ and luminosity $L$:
$F = \dfrac{L}{4\pi d^2}$.

For Vega:
\[
F_V = \frac{L_V}{4\pi\,d_V^2}
\]
For Supernova:
\[
F_{SN} = \frac{L_{SN}}{4\pi\,d_{SN}^2}
\]
\begin{align*}
F_{SN} = 1.6\times10^{-7}F_V
\quad\Longrightarrow\quad
\frac{L_{SN}}{4\pi\,d_{SN}^2}
  &= \left(1.6\times10^{-7}\right)\frac{L_V}{4\pi\,d_V^2}\\
L_{SN}\,d_V^2 &= \left(1.6\times10^{-7}\right)L_V\,d_{SN}^2
\end{align*}
\hfill(25\%)
\begin{align*}
d_{SN} &= d_V\left[\frac{1}{1.6\times10^{-7}}\,
  \frac{L_{SN}}{L_V}\right]^{1/2}
= 7.76\left[\frac{1}{1.6\times10^{-7}}\,
  \frac{5.8\times10^{9}}{130}\right]^{1/2}
= 129\,581\,812.59\ \mathrm{pc}\\
&= 1.3\times10^{8}\ \mathrm{pc}
\end{align*}
\hfill(35\%)

If $t_H$ is the Hubble time, then the relation between redshift, distance, and
Hubble time is:
\[
t_H = \frac{d_{SN}}{cz}
= \left[\frac{1.3\times10^{8}\ \mathrm{pc}}
  {(1\ \mathrm{lyr/yr})(0.03)}\right]
  \left(\frac{3.26\ \mathrm{lyr}}{1\ \mathrm{pc}}\right)
= 1.41\times10^{10}\ \mathrm{years}
\]
\hfill(40\%)

\solhead{13.3}{}
% source id: 2008_TH_S5
Universe change in size and redshift relation:
\[
1+z = \frac{\mathrm{size}_{\mathrm{now}}}{\mathrm{size}_{\mathrm{at}\,z}}
\quad\Longrightarrow\quad
\mathrm{size}_{\mathrm{at}\,z} = \frac{1}{1+z}\,\mathrm{size}_{\mathrm{now}}
= \frac{1}{1+1100}\,\mathrm{size}_{\mathrm{now}}
= 0.00091\,\mathrm{size}_{\mathrm{now}}
\]
so that the size of the Universe at the time of CMB emission was $0.00091$ of
its $\mathrm{size}_{\mathrm{now}}$.
\hfill(50\%)

If
\[
A = \frac{V_{\mathrm{CMB\ at\ emission\ time}}}{V_{\mathrm{Now}}}
= (0.00091)^3 = 7.5\times10^{-10},
\]
then the density of the Dark Matter at CMB emission time
\[
\rho_{\mathrm{DM\ at\ CMB}} = \frac{\rho_{\mathrm{DM}}}{A}
= \frac{2.4\times10^{-30}}{7.5\times10^{-10}}
= 3.2\times10^{-21}\ \mathrm{g/cm^3}.
\]
\hfill(25\%)

The ratio between the density of Dark Matter to the density of Dark Energy at
the time CMB was emitted
\[
= \frac{\rho_{\mathrm{DM\ at\ CMB}}}{\rho_{\Lambda}}
= \frac{3.2\times10^{-21}}{6.7\times10^{-30}} = 4.8\times10^{8}
\]
\hfill(25\%)

\solhead{13.4}{}
% source id: 2009_TH_P11
In a flat universe
\begin{equation*}
\rho = \rho_c = \frac{3H^2}{8\pi G} \tag*{(4 points)}
\end{equation*}
\begin{equation*}
H = 75~\mathrm{km\,s^{-1}\,Mpc^{-1}} = 2.4\times10^{-18}~\mathrm{s^{-1}}
\tag*{(2 points)}
\end{equation*}
\begin{equation*}
\rho_c = 1.1\times10^{-26}~\mathrm{kg\,m^{-3}} \tag*{(2 points)}
\end{equation*}
\begin{equation*}
n_\nu = \frac{0.25\,\rho_c}{10^{-5}\,m_e} = 3\times10^{8}~\mathrm{m^{-3}}
\tag*{(2 points)}
\end{equation*}

\solhead{13.5}{}
% source id: 2010_TH_S3
Recession velocity of the QSO is
\begin{equation*}
\frac{v}{c} = \frac{(z+1)^2 - 1}{(z+1)^2 + 1} = 0.18 \tag*{4}
\end{equation*}
According to the Hubble's law,
\begin{equation*}
v = H_0 D \tag*{2}
\end{equation*}
The distance of the QSO is
\begin{equation*}
D = v/H_0 = 0.18c/72 = 750~\mathrm{Mpc}, \tag*{4}
\end{equation*}

Remarks: if the student calculate the distance using cosmological formula and
arrive at the answer $D = 735~\mathrm{Mpc}$, assuming $\Omega_0 = 1.0$ will
get the full mark.

\solhead{13.6}{}
% source id: 2011_TH_S11
Within assumed cosmological model the following relation is valid:
\[
(1+z) = \frac{R_0}{R(z)} = \left(\frac{t_0}{t(z)}\right)^{2/3},
\]
where $R_0$ and $R(z)$ denote the current value of the scale factor and that
for the redshift equal to $z$ respectively, while $t_0$ and $t(z)$ stands for
the time that passed from the Big Bang. \hfill[5]

For $z=6.03$ the appropriate time $t(z) = 7.35\cdot10^{8}$ years. Subtracting
from $t(z)$ the interval equal to 560 My and 600 My one obtains
$t_1 = 1.75\cdot10^{8}$\,y and $t_2 = 1.35\cdot10^{8}$\,y \hfill[3]

what is equivalent to $z_1 = 17.3$ and $z_2 = 20.8$. \hfill[2]

\solhead{13.7}{}
% source id: 2011_TH_S15
Correct choice of the wavelength: \hfill[2]

If the nature of the spectrum does not change, it is sufficient to perform the
calculation for one wavelength. The simplest choice is for the wavelength of
the maximum (Wien's law)
\[
\lambda = b/T
\]

Understanding definition of redshift \hfill[2]

From this
\[
z = (\lambda_a - \lambda_e)/\lambda_e
\]
where $\lambda_a$ is the wavelength received at the Earth and $\lambda_e$ is
the emitted wavelength, thus $\lambda_e = \lambda_a/(z+1)$.

Substituting into Wien's law, we find
\[
T = b\cdot(z+1)/\lambda_a
\]
\hfill[4]

Substituting the current value of the temperature we find the result \hfill[2]

\solhead{13.8}{}
% source id: 2012_TH_S10
As we have a differential equation that should be integrated, we expect that
the student can solve it in two ways: using an approximation (Way 1) or
calculus (Way 2). \hfill(2 pt finding final scale factor)

\medskip
\noindent\textit{Way 1.}
Using the definition of Hubble constant:
$\dfrac{da}{dt} = Ha$, where $a$ is the scale factor and using the
approximation $\dfrac{da}{dt} \approx \dfrac{\Delta a}{\Delta t}$, and
supposing that the variation of the scale factor and the Hubble constant is
small compared to their values, and using that $aT = \mathrm{const}$:
\begin{equation*}
\frac{\Delta a}{\Delta t} = H_i a_i
\;\Rightarrow\;
\frac{\Delta T}{\Delta t} = -H_i T_i
\;\Rightarrow\;
\Delta t = -\frac{1}{H_i}\,\frac{\Delta T}{T_i} = 5.0\cdot10^{8}~years
\tag*{(7 pt + 1 pt)}
\end{equation*}

\medskip
\noindent\textit{Way 2.}
As we have a matter dominated Universe ($\rho \propto a^{-3}$), without dark
energy, and a density parameter $\Omega_0 = 1$, we have:
\begin{equation*}
\Omega_0 = \frac{8\pi G\rho}{3H^2} = \frac{8\pi G}{3H^2}\,\frac{\zeta}{a^3} = 1
\tag*{(4 pt)}
\end{equation*}
From the Hubble law, $\dfrac{da}{dt} = Ha$:
\[
\left(\frac{da}{dt}\right)^{2} = \frac{8\pi G}{3}\,\frac{\zeta}{a}
\;\Rightarrow\;
\frac{da}{dt} = \sqrt{\frac{8\pi G}{3}\,\frac{\zeta}{a}}
\]
Integrating this equation: \hfill(3 pt)
\[
\frac{2}{3}\left[a_f^{3/2} - a_i^{3/2}\right]
= \sqrt{\frac{8\pi G}{3}\,\zeta}\,\Delta t
\;\Rightarrow\;
\Delta t
= \frac{2}{3}\left[\left(\frac{a_f}{a_i}\right)^{3/2} - 1\right]
  \frac{1}{\sqrt{\dfrac{8\pi G}{3}\,\dfrac{\zeta}{a_i^{3}}}}
= \frac{2}{3}\,\frac{1}{H_i}
  \left[\left(\frac{a_f}{a_i}\right)^{3/2} - 1\right]
\]
Using the fact that $aT = \mathrm{const}$:
\[
a_i T_i = a_f T_f
\;\Rightarrow\;
\frac{a_f}{a_i} = \frac{T_i}{T_f} = \frac{2.73}{2.63} = 1.038
\]
So the time is:
\begin{equation*}
\Delta t = 5.2\cdot10^{8}~years \tag*{(1 pt)}
\end{equation*}
Obs: (7 pt for using correctly that $a \propto t^{2/3}$)

\solhead{13.9}{}
% source id: 2013_TH_L1
\begin{description}
\item[(1)] The critical density $\rho_c = \dfrac{3H_0^2}{8\pi G}$
\hfill(9 points)

If the matter density parameter
$\Omega_m = \dfrac{\rho_m}{\rho_c}$ is 32\%, thus
$\rho_m = 0.32\,\dfrac{3H_0^2}{8\pi G}$.
\hfill(3 points)

From the latest estimate of
$H_0 = 67.8\ \dfrac{\mathrm{km\,s^{-1}}}{Mpc}$ we obtain
$\rho_m = 8.6\times10^{-27}$ kg m$^{-3}$
% sic: 8.6e-27 kg/m^3 is the critical density for H_0 = 67.8; the matter
% density with the 0.32 factor would be 2.8e-27 kg/m^3.
\hfill(4 points)

\item[(2)] The escape velocity is
$\upsilon_{esc} = \sqrt{\dfrac{2GM}{d}}$.
\hfill(4 points)

By replacing $M = \rho_c\,\dfrac{4\pi}{3}d^3$
\hfill(2 points)

we obtain, for the escape velocity within $d = 100$ Mpc
\[
\upsilon_{esc}
= \sqrt{\frac{8\pi}{3}G\times0.32\times\frac{3H_0^2}{8\pi G}}
  \times 100\,Mpc
= 3835\ \mathrm{km\,s^{-1}}
\]
\hfill(8 points)

\item[(3)] If a galaxy is orbiting around the centre of our galaxy, its
velocity is $\dfrac{1}{\sqrt{2}}$ of its escape velocity. Thus
\[
\omega = \frac{v}{d} = \frac{v_{esc}/\sqrt{2}}{d}
= \frac{H_0 d\sqrt{\Omega_m/2}}{d}
= \frac{\sqrt{0.32}\,(67.8\times10^{3}\,m/s)/(3.09\times10^{22}\,m)}
  {\sqrt{2}}
= 8.8\times10^{-19}\ rad/s
\]
\hfill(12 Points)

This is $1.8\times10^{-13}\ arc\sec/s$ and it does not depend on the
distance $d$.

\item[(4)] Therefore we will never be able to resolve them and the answer is
``No''.
\hfill(8 Points)
\end{description}

\solhead{13.10}{}
% source id: 2013_TH_L3
\begin{description}
\item[(1)] Using the virial theorem for our isolated, spherical system of $N$
galaxies of mass $m$, each, we get
\[
-2\langle K\rangle = \langle U\rangle
\;\Rightarrow\;
-\frac{2}{N}\sum_{1}^{N}\frac{1}{2}m_i u_i^2 = \;<U>
\;\Rightarrow\;
-\frac{m}{N}\sum_{1}^{N}u_i^2 = \frac{U}{N}
\]
\hfill(3 point)

where
$\dfrac{1}{N}\displaystyle\sum_{1}^{N}u_i^2 \approx \langle u^2\rangle
= \langle u_r^{\,2}\rangle + \langle u_\theta^{\,2}\rangle
+ \langle u_\phi^{\,2}\rangle$,
where $u_r$, $u_\theta$ and $u_\phi$ are the radial velocity and the two
perpendicular velocities on the plane of the sky of the members of the
cluster.
\hfill(6 points)

Assuming that
$\langle u_r^{\,2}\rangle \sim \langle u_\theta^{\,2}\rangle
\sim \langle u_\phi^{\,2}\rangle$, we have
\[
-3m\langle u_r^2\rangle = U/N
= \Bigl(-\frac{3}{5}\frac{GM^2}{R}\Bigr)/N
= -\frac{3}{5}\frac{GMm}{R}
\;\Rightarrow\;
M \approx \frac{5R\langle u_r^2\rangle}{G}
\]
\hfill(13 points)

(where we used the gravitational potential of a spherical homogeneous mass
$M$ enclosed within radius $R$)

Alternatively the student can give a rougher order of magnitude estimate
\[
-2\langle K\rangle = \langle U\rangle
\;\Rightarrow\;
-2.\frac{1}{2}\langle u_r^2\rangle \approx -\frac{GM}{R}
\quad\text{etc}
\]
If they do not use the exact formula:
\hfill(11 Points)

\item[(2)] From the result of the previous question we have
$M \approx \dfrac{5R\langle u_r^2\rangle}{G}$, where
\[
\sqrt{\langle u_r^2\rangle} = \sigma = 1000\ km\,s^{-1}
\]
\hfill(8 points)

The angular diameter of the cluster is $\varphi = 4^\circ$ at a distance of
$d = 90$ Mpc. Therefore the diameter D of the cluster in Mpc is calculated
from:
\[
\tan\phi = \frac{D}{d}
\;\Rightarrow\;
\phi(rad) \approx \frac{D}{d}
\;\Rightarrow\;
D \approx 4\times\frac{\pi}{180}\times90\,\mathrm{Mpc} \approx 6.3\,Mpc
\;\Rightarrow\;
R \approx 3\,Mpc
\]
\hfill(11 points)

Therefore
$M \approx \dfrac{5R\langle u_r^2\rangle}{G} = 3.6\times10^{15}\,\mathrm{M}_\odot$
\hfill(3 point)

\item[(3)]
\[
\frac{M_{virial}}{L_{galaxies}}
= \frac{3.6\times10^{15}\,\mathrm{M}_\odot}{5\times10^{12}\,L_\odot}
= 720\,\frac{\mathrm{M}_\odot}{L_\odot}
\]
This is obviously much larger from
$\dfrac{\mathrm{M}_\odot}{L_\odot} = 1$ which is found from the visible mass
of the cluster. Thus $\dfrac{M}{M_{virial}} = \dfrac{1}{720}$.
\hfill(6 points)
\end{description}

\solhead{13.11}{}
% source id: 2013_TH_S15
The recession velocity is calculated by either the classical relation,
$v_c = z\times c$, or the relativistic relation,
\[
v_r = \frac{(1+z)^2 - 1}{(1+z)^2 + 1}\,c.
\]
The calculations for the three galaxies are summarized in the following
Table:

\begin{center}
\begin{tabular}{|c|c|c|c|c|}
\hline
Galaxy & $v_c$ (km/s) & $V_a$ (km/s) & $v_r$ (km/s) & $v_r/c \times 100$\\
\hline
3C279 & 160800 & 128750 & 121390 & 40\%\\
\hline
3C245 & 308700 & 212260 & 182740 & 61\%\\
\hline
4C41.17 & 1140000 & 470580 & 275040 & 92\%\\
\hline
\end{tabular}
\end{center}

\begin{description}
\item[(1)] Table 1, Columns 2,3,4 \hfill(5 Points)
\item[(2)] Column 5 \hfill(3 Points)
\item[(3)] If the student answers ``Classical'' \hfill(0 Points)\\
  If the student answers ``Special Relativity'' \hfill(1 Point)\\
  If the student answers ``Approximate formula'' \hfill(2 Points)
\end{description}

\solhead{13.12}{}
% source id: 2013_TH_S5
\begin{description}
\item[(a)]
\[
z = \frac{\lambda - \lambda_o}{\lambda_o}
= \frac{662.9 - 656.3}{656.3} \approx 0.01.
\]
\hfill(3 Point)

This is small enough that we can use the classical equation for the expansion
of the Universe.
\[
H_o = \frac{cz}{r}
= \frac{3\times10^{5}\ \mathrm{km\,s^{-1}} \times 10^{-2}}{41.67\ Mpc}
\quad\text{or}\quad
H_o = 72.4\ \frac{\mathrm{km\,s^{-1}}}{Mpc}
\]
\hfill(4 Points)

\item[(b)] $t_H \approx \dfrac{1}{H_o}$ or $t_H = 13.5$\,\textbf{Gyr}
\hfill(3 Point)
\end{description}

\solhead{13.13}{}
% source id: 2015_TH_S14
To decide how far one can see on average in a universe filled with spherical
objects of radius $R$, it is simplest to think of a long cylinder along the
line of sight. If an object is closer than $R$ to the line of sight, then the
line of sight intersects its surface. For a distance $\ell$, the cylindrical
volume which would contain such objects is $\pi R^2\ell$. If the density of
objects is $n$, then the volume that will on average contain one object is
defined by $\pi R^2\ell n = 1$ and the average distance to which we see
before our vision is blocked is
\[
\ell = \frac{1}{\pi R^2 n}
\]
\hfill(50 points)

The average radius $R$ in Mpc
\[
R = \frac{6.96\times10^{8}\ \mathrm{m}}
  {3.086\times10^{22}\ \mathrm{m\,Mpc^{-1}}}
= 2.3\times10^{-14}\ \mathrm{Mpc}
\]
\hfill(25 points)

Then
\[
\ell = \frac{1}
  {3.14\,(2.3\times10^{-14}\ \mathrm{Mpc})^2 (10^{9}\ \mathrm{Mpc^{-3}})}
= 6.3\times10^{17}\ \mathrm{Mpc}
\]
\hfill(25 points)

\solhead{13.14}{Early Universe}
% source id: 2016_TH_T3
\begin{description}
\item[(T3.1)]
\[
\frac{\rho_{\mathrm{m_0}}/\rho_{\mathrm{c}}}{\rho_{\mathrm{r_0}}/\rho_{\mathrm{c}}}
= \frac{\Omega_{\mathrm{m_0}}}{\Omega_{\mathrm{r_0}}}
= \frac{0.3}{10^{-4}} = 3000
\]
At $z_{\mathrm{e}}$, both matter density and radiation density were equal.
\begin{align*}
\rho_{\mathrm{r}} &= \rho_{\mathrm{m}}\\
\therefore\ \rho_{\mathrm{r_0}}(1+z_{\mathrm{e}})^4
  &= \rho_{\mathrm{m_0}}(1+z_{\mathrm{e}})^3\\
1+z_{\mathrm{e}} &= \frac{\rho_{\mathrm{m_0}}}{\rho_{\mathrm{r_0}}} = 3000
  \tag*{(2.0)}\\
\therefore\ z_{\mathrm{e}} &\simeq 3000 \tag*{(1.0)}
\end{align*}
\textbf{Only $z_{\mathrm{e}} = 2999$ and $z_{\mathrm{e}} = 3000$ are
acceptable answers.}

\item[(T3.2)] As the Universe behaves like an ideal black body, the radiation
density will be proportional to the fourth power of the temperature (Stefan's
law).
\begin{align*}
\left(\frac{T_{\mathrm{e}}}{T_0}\right)^4
  &= \frac{\rho_{\mathrm{r_e}}}{\rho_{\mathrm{r_0}}}
   = \frac{\rho_{\mathrm{r_0}}(1+z_{\mathrm{e}})^4}{\rho_{\mathrm{r_0}}}
  \tag*{(2.0)}\\
\left(\frac{T_{\mathrm{e}}}{2.732}\right)^4
  &= (1+z_{\mathrm{e}})^4 \tag*{(1.0)}\\
\frac{T_{\mathrm{e}}}{2.732} &= 1+z_{\mathrm{e}} = 3000\\
T_{\mathrm{e}} &= 3000\times2.732\\
T_{\mathrm{e}} &= 8200\ \mathrm{K} \tag*{(1.0)}
\end{align*}
\textbf{$8100 \le T_{\mathrm{e}} \le 8200$ gives 1.0;
$8200 < T_{\mathrm{e}} \le 9000$ gives 0.5; else 0.}

\item[(T3.3)] Wien's law:
\begin{align*}
\lambda_{\max} &= \frac{0.002898\ \mathrm{m\,K}}{T_e}\\
  &= \frac{0.002898}{8200}\ \mathrm{m} = 354\ \mathrm{nm} \tag*{(1.0)}\\
E_\nu &= \frac{hc}{\lambda_{\max}}\\
  &= \frac{6.62\times10^{-34}\times3\times10^{8}}{354\times10^{-9}}\
     \mathrm{J}
   = \frac{5.62\times10^{-19}}{1.602\times10^{-19}}\ \mathrm{eV}
  \tag*{(1.0)}\\
E_\nu &= 3.5\ \mathrm{eV} \tag*{(1.0)}
\end{align*}
Alternative solution:
\begin{align*}
E_\nu &= k_{\mathrm{B}}T_{\mathrm{e}}\\
  &= 1.38\times10^{-23}\times8200\ \mathrm{J}
   = \frac{1.13\times10^{-19}}{1.602\times10^{-19}}\ \mathrm{eV}
  \tag*{(2.0)}\\
E_\nu &= 0.71\ \mathrm{eV} \tag*{(1.0)}
\end{align*}
\textbf{Use of either Wien's law or $E = k_{\mathrm{B}}T$ gets full credit.
$E_\nu = 3k_{\mathrm{B}}T/2$ or similar gets no credit. Answers with
$E_\nu = 3k_{\mathrm{B}}T$ or $E_\nu = 2.7k_{\mathrm{B}}T$ also get full
credit.}
\end{description}

\solhead{13.15}{Super Luminal Galaxies}
% source id: 2018_TH_T1
The statements are, in order:

\begin{description}
\item[(a)] True
\item[(b)] False
\item[(c)] True
\item[(d)] False
\item[(e)] False
\end{description}

\hfill(2 points each)

\solhead{13.16}{Thermal History of the Universe}
% source id: 2018_TH_T11
\begin{description}
\item[(a)] One can argue this in different ways. One can note that the
  Hubble parameter is often called Hubble constant, but actually it's not a
  constant. The dimensions of Hubble parameter is the inverse of time
  $[T^{-1}]$. Alternatively, one can simply look at unit of $H_0$ in the
  table of constants and conclude the same. \hfill(1 point)

  It's natural to define a timescale as the reciprocal of Hubble
  parameter: $t_H = \frac{1}{H}$, and this is the Hubble time which is a
  characteristic timescale of Universe expansion. \hfill(2 points)

  Present-day Hubble time $t_{H0} = 14.46$\,Gyr. \hfill(2 points)

\item[(b)] From Friedmann Equation, the critical density is defined in the
  way of
  $\left(\frac{\dot{a}}{a}\right)^2 = \frac{8\pi G}{3}\rho_c$, thus
  $\rho_c = \frac{3H^2}{8\pi G}$. \hfill(2 points)

  Substitute $H_0 = 2.19\times 10^{-18}\,\mathrm{s^{-1}}$ into above, we
  have $\rho_{c0} = 8.59\times 10^{-27}\ kg/m^3$ \hfill(3 points)

\item[(c)] $\Omega_{\mathrm{m}} + \Omega_{\mathrm{r}} + \Omega_\Lambda
  + \Omega_k = 1$
  \begin{align*}
  \frac{\rho_m}{\rho_c} + \frac{\rho_r}{\rho_c}
    + \Omega_\Lambda + \Omega_k &= 1\\
  \frac{8\pi G}{3H^2}(\rho_m + \rho_r)
    + \Omega_\Lambda + \Omega_k &= 1\\
  \frac{8\pi G}{3}(\rho_m + \rho_r)
    + H^2\Omega_\Lambda + H^2\Omega_k &= H^2
  \end{align*}
  \hfill(2 points)

  Comparing this with the Friedmann equation
  \[ H^2 = \frac{8\pi G}{3}(\rho_m + \rho_r)
     + \frac{\Lambda c^2}{3} - \frac{kc^2}{a^2}. \]
  \[ \Omega_\Lambda = \frac{\Lambda c^2}{3H^2}
     \quad\text{and}\quad
     \Omega_k = -\frac{kc^2}{a^2 H^2} \]
  \hfill(4 points)

\item[(d)] Pressure exerted by the radiation is
  $p = \frac{1}{3}\rho_r c^2$

  Thus, the Fluid Equation becomes,
  $\dot{\rho_r} + 3\frac{\dot{a}}{a}
   \left(\rho_r + \frac{1}{3}\rho_r\right) = 0$ \hfill(2 points)

  $\frac{\dot{\rho_r}}{\rho_r} = -4\frac{\dot{a}}{a}$ implying

  Thus $\rho_r \propto a^{-4}$. \hfill(3 points)

  And we know $a = \frac{1}{1+z}$, hence $\rho_r \propto (1+z)^4$.
  \hfill(2 points)

\item[(e)] In the fluid equation,
  $\dot{\rho} + 3\frac{\dot{a}}{a}(\rho + q\rho) = 0$.

  But the cosmological constant does not evolve. So
  $\dot{\rho_\Lambda} = 0$

  $\frac{3\dot{a}\rho_\Lambda}{a}(1 + q) = 0$ which means $q = -1$.
  \hfill(4 points)

\item[(f)] From the units, one can figure out that $\hbar$ has the
  dimensions of $[ML^2T^{-1}]$.

  $G$ has the dimensions of $[M^{-1}L^3T^{-2}]$.

  And speed of light $c$ has dimensions of $[LT^{-1}]$. \hfill(5 points)

  We assume that Planck time is given by $\hbar^x c^y G^z$. In order to
  make it a time dimensional quantity, the following equation must hold:
  $T = (ML^2T^{-2})^x(LT^{-1})^y(M^{-1}L^3T^{-2})^z$.
  % sic: source writes (ML^2 T^{-2})^x for hbar despite stating
  % [ML^2 T^{-1}] above; the exponent bookkeeping below uses T^{-1}
  $T = M^{x-z}L^{2x+y+3z}T^{-x-y-2z}$. \hfill(3 points)

  From $x - z = 0$, $2x + y + 3z = 0$, $-x - y - 2z = 1$, we get
  $x = \frac{1}{2}$, $y = \frac{-5}{2}$, $z = \frac{1}{2}$. Thus, the
  Planck time is approximately $t_p \approx \sqrt{\hbar G/c^5}$.
  \hfill(3 points)

  Substitute the numerical values into it and get:
  $t_P = 5.4\times 10^{-44}\,s$. \hfill(2 points)

\item[(g)] Schwarzschild radius of a black hole equals to $2l_P$:
  \[ r_s = \frac{2GM_P}{c^2} = 2ct_P = 2c\sqrt{\hbar G/c^5}, \]
  \hfill(3 points)
  \[ M_P = \sqrt{\hbar c/G}. \]
  \hfill(2 points)
  \[ M_P c^2 = \sqrt{\hbar c^5/G} = 1.22\times 10^{19}\,GeV. \]
  \hfill(2 points)

\item[(h)] Given the energy density, and the result of problem \textbf{d},
  we find that $\varepsilon_r \propto T^4$,
  $\varepsilon_r = \rho_r c^2$, $\rho_r \propto (1+z)^4$.

  Thus $T \propto 1 + z$, $\frac{T}{1+z} = const$. \hfill(4 points)

\item[(i)] The matter-radiation equality means
  $\Omega_r(z_{eq}) = \Omega_m(z_{eq})$. We need to derive the behavior of
  this two density parameters. Obviously
  $\Omega_m(z) = \Omega_{m0}(1+z)^3$ and
  $\Omega_r(z) = \Omega_{r0}(1+z)^4$. \hfill(4 points)

  Thus $1 + z_{eq} = \frac{\Omega_{m0}}{\Omega_{r0}}$. $\Omega_{r0}$ is
  the density parameter of CMB radiation and can be calculated.
  \hfill(1 point)

  By the definition of density parameter of photon radiation (footnote
  gamma: $\gamma$,
  \[ \Omega_{\gamma 0} = \frac{\rho_{\gamma 0}}{\rho_{c0}}
     = \frac{\varepsilon_{\gamma 0}}{c^2}\frac{8\pi G}{3H^2}
     = \frac{\varepsilon_{\gamma 0}}{c^2}
       \frac{8\pi G}{3H_{100}^2 h^2}, \]
  where $H_{100} = 100\ km/s/Mpc$. \hfill(4 points)

  The energy density of blackbody photon radiation is
  $\varepsilon_\gamma = \frac{\pi^2}{15\hbar^3 c^3}(k_B T)^4$, thus
  \[ \Omega_{\gamma 0}h^2 = \frac{\pi^2}{15\hbar^3 c^5}\,k_B^4\,
     \frac{8\pi G}{3H_{100}^2}\,T_0^4 = 2.47\times 10^{-5}. \]
  \hfill(4 points)

  But there is also an additional contribution from neutrinos of about
  68\%, so the total density parameter of radiation is
  $\Omega_{r0}h^2 = \Omega_{\gamma 0}h^2\times 1.68
   = 4.15\times 10^{-5}$. \hfill(2 points)

  Hence
  $1 + z_{eq} = \frac{\Omega_{m0}}{\Omega_{r0}}
   = 2.4\times 10^{4}\,\Omega_{m0}h^2$. \hfill(1 point)

\item[(j)] At that time, the redshift is
  \[ 1 + z = \frac{1\,MeV}{k_B\times 2.73\,K} = 4.25\times 10^{9}. \]

  Thus the scale factor
  $a = \frac{1}{1+z} = 2.35\times 10^{-10}$. \hfill(2 points)

  Only considering radiation dominated Universe, Friedmann Equation can be
  written as
  \[ \left(\frac{\dot{a}}{a}\right)^2 = \frac{8\pi G}{3}\rho_r
     = \frac{8\pi G\rho_{r0}}{3a^4}, \]
  \hfill(2 points)
  \[ \frac{\dot{a}}{a} = \frac{1}{a^2}\sqrt{\frac{8\pi G\rho_{r0}}{3}} \]
  \[ a^2 H = \sqrt{\frac{8\pi G\Omega_{r0}}{3}\rho_{c0}}
     = \sqrt{H_0^2\Omega_{r0}} \]
  \hfill(1 points)

  Thus, we have,
  \[ t = \frac{1}{2H} = \frac{a^2}{2\sqrt{H_0^2\Omega_{r0}}}
     = 1.3\,s \sim 1\,s. \]

  Thus neutrino decoupled at about 1\,s after big bang. \hfill(3 points)
\end{description}

\solhead{13.17}{Distance to the Coma galaxy cluster}
% source id: 2019_TH_13
\begin{description}
\item[a)] The mean radial velocity of the cluster, from the data given in
  the table, is
  \[ \langle v_{\mathrm{r}}\rangle \approx 6731\ \mathrm{km\,s^{-1}} \]
  \hfill(5\,p)

  Using the Hubble law in the following form
  \[ v_{\mathrm{r}} = H_0 D, \]
  \hfill(2\,p)

  where $v_{\mathrm{r}}$ is the mean radial velocity in km/s, $D$ is the
  distance in Mpc, and $H_0 = 70\ \mathrm{km\,s^{-1}\,Mpc^{-1}}$ is the
  Hubble constant, we can get
  \[ D = \frac{v_{\mathrm{r}}}{H_0} \]
  With numerical values:
  \[ D = 96\ \mathrm{Mpc} \]
  \hfill(1\,p)

\item[b)] The angular diameter of the cluster is
  $\theta = 100' = 6000''$. Applying the definition of the parsec, $1''$
  angular distance from 1\,pc distance corresponds to 1\,au (Astronomical
  Unit, $1\,\mathrm{au} \approx 1.5\times 10^{8}$\,km), i.e.
  \[ d\,[\mathrm{AU}] = \theta\,['']\times D\,[\mathrm{pc}] \]
  \hfill(3\,p)

  Therefore, $\theta = 6000''$ from
  $D = 96\ \mathrm{Mpc} = 96\times 10^{6}$\,pc gives (with
  $1\,\mathrm{Mpc} = 3.086\times 10^{22}$\,m):
  \[ d = 5.77\times 10^{11}\,\mathrm{au} = 8.63\times 10^{22}\,\mathrm{m}
     = 2.80\ \mathrm{Mpc} \]
  \hfill(1\,p)

  Alternative solution: $100' = 1.67^\circ = 0.0291$\,rad. For such small
  angles $d \approx \theta D$. If $\theta$ is expressed in radians, then
  $d$ and $D$ have the same units. Thus, if $D = 96$\,Mpc, then
  \[ d = 0.0291\times 96\ \mathrm{Mpc} = 2.80\ \mathrm{Mpc} \]

\item[c)] In the virial theorem
  \[ \text{(13.1)}\quad \langle K\rangle
     = \frac{1}{N}\sum_{i=1}^{N}\frac{1}{2}\,m_i(v_i - v_{\mathrm{m}})^2 \]
  is the mean kinetic energy of $N$ galaxies. The $i^{\mathrm{th}}$ galaxy
  has mass $m_i$ and velocity $v_i$, and $v_{\mathrm{m}}$ is the mean
  velocity of the cluster.
  \[ \text{(13.2)}\quad \langle U\rangle
     = -\frac{3}{5}\,\frac{1}{N}\sum_{i=1}^{N}\frac{GM}{R}\,m_i \]
  is the mean gravitational potential energy of $N$ galaxies (having
  individual masses $m_i$ as above) filling in a sphere of $R$ radius,
  while $G$ is the gravitational constant and
  $M = \sum_{i=1}^{N} m_i$ is the total mass of the cluster.

  If the orbital velocity vectors are distributed randomly, then
  \[ \text{(13.3)}\quad
     \frac{1}{N}\sum_{i=1}^{N}(v_i - v_{\mathrm{m}})^2
     = \frac{1}{N}\sum_{i=1}^{N} 3(v_{i,\mathrm{r}} - v_{\mathrm{m,r}})^2
     = \frac{N-1}{N}\sum_{i=1}^{N}
       \frac{3(v_{i,\mathrm{r}} - v_{\mathrm{m,r}})^2}{N-1}
     \to 3\sigma_{\mathrm{r}}^2,\ \text{if}\ N\to\infty, \]
  where $v_{i,\mathrm{r}}$ is the radial velocity for the
  $i^{\mathrm{th}}$ galaxy, and $\sigma_{\mathrm{r}}$ is called the
  \emph{velocity dispersion} of the cluster.

  Combining Eq.~(13.1) and Eq.~(13.3) one can get:
  \[ \langle K\rangle = \frac{3}{2}\,m\sigma_{\mathrm{r}}^2 \]
  \hfill(2\,p)

  Inserting $\langle K\rangle$ into the equation of virial theorem gives:
  \[ -2\times\frac{3}{2}\,m\sigma_{\mathrm{r}}^2
     = -3m\sigma_{\mathrm{r}}^2 = \langle U\rangle \]
  \hfill(2\,p)

  Using Eq.~(13.2) to express $\langle U\rangle$ one can get
  \[ \langle U\rangle
     = -\frac{3}{5}\,m\,\frac{GM}{R}\,\frac{1}{N}\sum_{i=1}^{N} 1
     = -\frac{3}{5}\,m\,\frac{GM}{R}, \]
  \hfill(2\,p)

  because $\sum_{i=1}^{N} 1 = N$.

  Inserting this expression into the one above, we get:
  \[ -3m\sigma_{\mathrm{r}}^2 = -\frac{3}{5}\,m\,\frac{GM}{R} \]
  \hfill(2\,p)

  Dividing by $-3m$ both sides, and expressing $M$ we finally arrive at
  \[ M = \frac{5R\sigma_{\mathrm{r}}^2}{G}, \]
  \hfill(2\,p)

  which is the expression of the virial mass.

\item[d)] The radial velocity dispersion of the data given in the table is
  the same as the standard deviation of the data around the mean. The mean
  velocity, $\langle v_{\mathrm{r}}\rangle = 6731\ \mathrm{km\,s^{-1}}$,
  was already derived in part a). Thus (depending on how we define the
  standard deviation):
  \[ \sigma_{\mathrm{r},N}
     = \sqrt{\frac{1}{N}\sum_{i=1}^{N}
       \left(v_{\mathrm{r}} - \langle v_{\mathrm{r}}\rangle\right)^2}
     = 1184\ \mathrm{km\,s^{-1}} \]
  or
  \[ \sigma_{\mathrm{r},N-1}
     = \sqrt{\frac{1}{N-1}\sum_{i=1}^{N}
       \left(v_{\mathrm{r}} - \langle v_{\mathrm{r}}\rangle\right)^2}
     = 1214\ \mathrm{km\,s^{-1}} \]
  \hfill(10\,p)

  Inserting this value into the expression of the virial mass and using
  $R \approx 1.40$\,Mpc as the cluster radius (from b)), the result is:
  \begin{align*}
  M_N &= \frac{5\times(1.40\times 3.086\times 10^{22}\,\mathrm{m})
    \times(1.184\times 10^{6}\,\mathrm{m\,s^{-1}})^2}
    {6.67\times 10^{-11}\,\mathrm{m^3\,kg^{-1}\,s^{-2}}}\\
  M_N &= 4.53\times 10^{45}\,\mathrm{kg}
    \approx 2.28\times 10^{15}\,M_\odot
  \end{align*}
  or
  \begin{align*}
  M_{N-1} &= \frac{5\times(1.40\times 3.086\times 10^{22}\,\mathrm{m})
    \times(1.214\times 10^{6}\,\mathrm{m\,s^{-1}})^2}
    {6.67\times 10^{-11}\,\mathrm{m^3\,kg^{-1}\,s^{-2}}}\\
  M_{N-1} &= 4.77\times 10^{45}\,\mathrm{kg}
    \approx 2.40\times 10^{15}\,M_\odot
  \end{align*}
  \hfill(2\,p)

\item[e)] Using the virial mass from d) and the given total luminosity of
  the cluster one can get:
  \[ \frac{M_N}{L} = \frac{2.28\times 10^{15}\,M_\odot}
     {5\times 10^{12}\,L_\odot}
     = 455\ M_\odot/L_\odot \approx 460\ M_\odot/L_\odot \]
  or
  \[ \frac{M_{N-1}}{L} = \frac{2.40\times 10^{15}\,M_\odot}
     {5\times 10^{12}\,L_\odot}
     = 479\ M_\odot/L_\odot \approx 480\ M_\odot/L_\odot \]
  \hfill(1\,p)

  This mass-luminosity ratio is much higher than that of the individual
  galaxies which is typically in between 1 and 100. \hfill(1\,p)

\item[f)] (A) and (D) \hfill(4\,p)
\end{description}

\solhead{13.18}{Mirror}
% source id: 2020_TH_3
\begin{description}
\item[(a)] Let $t_1$ be the time taken by the light beam from the Solar
  System to the mirror, let $t_2$ be the time taken by the beam from the
  mirror to the Solar System and $T$ the total time to go back and forth.
  As a first order approximation, we will take distance travelled by the
  photon in each part as an average of the initial and final distance.
  Therefore, equating the kinematics of the situation, we have:
  \begin{align*}
  S_1 &= \frac{D + De^{H_0 t_1}}{2}
       = \frac{D\left(1 + e^{H_0 t_1}\right)}{2} = ct_1\\
  S_2 &= \frac{S_1 + S_1 e^{H_0 t_2}}{2}
       = \frac{S_1\left(1 + e^{H_0 t_2}\right)}{2} = ct_2
  \end{align*}
  \hfill(1.0)
  \begin{align*}
  &= \frac{D}{4}\left(1 + e^{H_0 t_1}\right)\left(1 + e^{H_0 t_2}\right)\\
  &\approx \frac{D}{4}\left(2 + H_0 t_1\right)\left(2 + H_0 t_2\right)\\
  &\approx \frac{D}{4}\left[4 + 2H_0\left(t_1 + t_2\right)\right]\\
  S_2 &= D\left(1 + \frac{1}{2}H_0 T\right)
  \end{align*}
  \hfill(1.5)

  From the first equation, we also find:
  \begin{align*}
  S_1 = ct_1 &= \frac{D\left(1 + e^{H_0 t_1}\right)}{2}\\
  ct_1 &= D\left(1 + \frac{1}{2}H_0 t_1\right)\\
  \therefore\ t_1 &= \frac{D}{c - \frac{1}{2}DH_0}\\
  S_1 = ct_1 &= \frac{2Dc}{2c - DH_0}
  \end{align*}
  \hfill(1.5)
  \[ \text{similarly, } S_2 = \frac{2S_1 c}{2c - S_1 H_0}. \]
  \hfill(0.5)

  Joining the expressions found, we obtain:
  \begin{align*}
  S_2 = D\left(1 + \frac{1}{2}H_0 T\right)
    &= \frac{2S_1 c}{2c - S_1 H_0}\\
  D\left(1 + \frac{1}{2}H_0 T\right)
    &= \frac{\frac{4Dc^2}{2c - DH_0}}
            {2c - \frac{2Dc}{2c - DH_0}\cdot H_0}\\
  \left(1 + \frac{1}{2}H_0 T\right)
    &= \frac{2c}{2c - DH_0 - DH_0} = \frac{c}{c - DH_0}\\
  c &= \left(1 + \frac{1}{2}H_0 T\right)(c - DH_0)\\
    &= c - DH_0 + \frac{1}{2}cH_0 T - \frac{1}{2}DH_0^2 T\\
  0 &= \frac{H_0}{2}\left(cT - 2D - DH_0 T\right)\\
  H_0 &= \frac{cT - 2D}{DT}.
  \end{align*}
  \hfill(2.5)

  \textbf{Alternative solution}\\
  First note that the time taken for the laser beam to travel to the mirror
  and back again is equal to the time taken for the laser beam to travel a
  distance of $2D$ (measured at $t = 0$) in a straight line. \hfill(2.0)\\
  The co-moving distance the beam has to cover is the same in both
  scenarios. We then equate the distance travelled by the light with the
  amount by which the space has expanded in that time.

  In time $t$, the space has expanded by a factor of $\exp(H_0 t)$, to
  first order we can linearise this to $1 + H_0 t$. This means that the
  beam, on average, travels through space that's stretched out by a factor
  of $1 + H_0 T/2$. \hfill(2.5)\\
  Thus,
  \[ cT = 2D\left(1 + \frac{H_0 T}{2}\right) \]
  and so
  \[ H_0 = \frac{cT - 2D}{DT}. \]
  \hfill(2.5)

  \textbf{Better approximation (not required in the exam)}\\
  In reality, the light does not travel a distance of $ct$ in this
  expanding universe. To find the exact expression, consider a time
  differential $\mathrm{d}t$ time $t$ after the beam was emitted. In this
  small time interval, the beam travels a distance of $c\,\mathrm{d}t$.
  Because the space is stretched out, the travelled distance corresponds to
  a smaller segment of space at $t = 0$, smaller by a factor of
  $\exp(H_0 t)$. The distance spanned at $t = 0$ is then
  \[ \mathrm{d}r = \exp(-H_0 t)\,c\,\mathrm{d}t. \]
  We integrate this from $t = 0$ to $t = T$:
  \begin{align*}
  \int_0^{2D}\mathrm{d}r &= c\int_0^{T}\exp(-H_0 t)\,\mathrm{d}t,\\
  2D &= \frac{c}{H_0}\left(1 - \exp(-H_0 T)\right).
  \end{align*}
  The result so far is accurate within the constraints of the model, but
  it's not analytically solvable for $H_0$. To get an estimate, we can
  approximate the right hand side to
  \[ 2D = \frac{c}{H_0}\left(1 - 1 + H_0 T - \frac{(H_0 T)^2}{2}\right)
     = cT - \frac{cH_0 T^2}{2}. \]
  Expressing $H_0$, we get
  \[ H_0 = \frac{2}{cT^2}\left(cT - 2D\right). \]
  The difference between this answer and the initial estimate is $2D/cT$
  which is almost unitary.

\item[(b)] From the $H_0$ expression found in the previous item, we find
  the travel time
  \begin{align*}
  T &= \frac{2D}{c - DH_0}
  \end{align*}
  \hfill(1.0)
  \begin{align*}
  &= \left(\frac{c}{2D} - \frac{H_0}{2}\right)^{-1}
   \approx \left(\frac{c}{2D}\right)^{-1} = \frac{2D}{c}
  \end{align*}
  \hfill(1.0)
  \begin{align*}
  T &= \frac{2\times 8\times 3.086\times 10^{16}}{3\times 10^{8}}\\
    &= 1.65\times 10^{9}\,\mathrm{s}\\
  T &\approx 52.2\,\mathrm{yr}.
  \end{align*}
  \hfill(1.0)
\end{description}

\solhead{13.19}{Cosmic String}
% source id: 2021_TH_Q15
\textbf{PART A:}

\begin{description}
\item[A.1] In cylindrical coordinates, the gravitational field is given by
  \[ \vec{g} = \vec{g}(r) = g(r)\,\hat{r} = g\,\hat{r}, \]
  \hfill[1 point]

  where, in view of the axial symmetry, $g$ can be computed from the Gauss
  law for the gravitational field:
  \[ \vec{g}\bullet\vec{A} = -\,4\pi G M_{in}, \]
  \hfill[1 point]

  where $M_{in}$ is the mass enclosed by the surface $A$, as shown in the
  figure below:
  \ioaafig{2021_TH_Q15-s1.pdf}

  Taking the surface to be a cylinder of length $L$ and radius $r$, from
  the above equation:
  \[ 2\pi r L g = -\,4\pi G M_{in} \]
  \[ g = -\frac{2GM_{in}}{rL} \]
  \hfill[1 point]

  There are two cases (regions):
  \begin{itemize}
  \item For $r > r_0$, one has $M_{in} = \mu L$ and so \hfill[0.5 points]
    \[ g = -\frac{2G\mu}{r} \]
    \hfill[1 point]
  \item For $r < r_0$, one has $M = \frac{r^2}{r_0^2}\mu L$ and so
    \hfill[0.5 points]
    \[ g = -\frac{2G\mu}{r_0}\left(\frac{r}{r_0}\right) \]
    \hfill[1 point]
  \end{itemize}

\item[A.2] Writing $\vec{g}(r_0)$, in terms of $G$, $\mu$ and $r_0$:
  \begin{align*}
  \vec{g}(r_0) &= -\frac{2G\mu}{r_0}\hat{r}\\
  g_0 \equiv \left|\vec{g}(r_0)\right| &= \frac{2G\mu}{r_0}
  \end{align*}
  \hfill[1 point]

\item[A.3] The following points should be included in the sketch:
  \begin{itemize}
  \item $(0, 0)$ \hfill[0.5 points]
  \item $(1, -1)$ \hfill[0.5 points]
  \end{itemize}
  Additionally,
  \begin{itemize}
  \item the sketch is linear between $(0,0)$ and $(1,-1)$
    \hfill[0.8 points]
  \item the sketch is concave for $r > r_0$,
    $\left(\sim\frac{1}{r}\right)$ \hfill[0.8 points]
  \item the sketch is drawn below x-axis. \hfill[0.4 points]
  \end{itemize}
  \ioaafig{2021_TH_Q15-s2.pdf}

\item[A.4] The centripetal acceleration equals the gravitational
  acceleration, thus
  \[ \frac{v^2}{R} = \frac{2G\mu}{R} \]
  \hfill[1 point]
  \[ \frac{2\pi R}{\tau} = \sqrt{2G\mu} \]
  \hfill[1 point]
  \[ R = \frac{\sqrt{2G\mu}}{2\pi}\tau \]
  \[ A = \frac{\sqrt{2G\mu}}{2\pi} \]
  \hfill[1 point]
  \[ \alpha = 1 \]
  \hfill[1 point]

\item[A.5] The potential energy of the particle is given by
  \[ U(r) = -\int_b^r \vec{F}(r)\cdot\vec{dr} \]
  \[ = m\int_b^r g(r)\,dr \]
  \hfill[1 point]
  \[ = 2Gm\mu\int_b^r \frac{1}{r}\,dr \]
  \hfill[1 point]
  \[ = 2Gm\mu\,\ln\!\left(\frac{r}{b}\right) \]
  \hfill[1 point]

  Where $b$ is a constant that sets the 0 of $U$.

  Note 1: Usually, $U = 0$ is set at $b = \infty$. For a cosmic string,
  this is not possible. Instead, we can choose $U = 0$, for example, at
  the string's surface $b = r_0$ \emph{(this choice is not relevant!)}.
  Thus
  \[ U = 2Gm\mu\,\ln\!\left(\frac{r}{r_0}\right) \]

  Note 2: For completeness, notice that the potential inside the string,
  $r < r_0$ is
  \[ U = 2Gm\mu\left(\frac{r}{r_0}\right)^2 - U_0, \]
  where
  $U_0 = 2Gm\mu\left(\frac{r}{r_0}\right)^2
   - 2\,Gm\mu\,\ln\!\left(\frac{r}{r_0}\right)$
  is a constant (again, not relevant!) that can be chosen such that $U$
  is continuous at the surface of the string. However, for the solution
  of the problem, it is not necessary to show this.

\item[A.6] The total energy of the particle is conserved, thus:
  \[ \frac{1}{2}mv^2 + 2Gm\mu\,\ln\!\left(\frac{r}{b}\right)
     = 2Gm\mu\,\ln\!\left(\frac{R_{max}}{b}\right) \]
  \hfill[2 point]

  Solving for $R_{max}$:
  \[ R_{max} = Re^{\frac{v^2}{4G\mu}} \]
  \hfill[2 point]

  Notice that the answer does not depend on $b$.

\item[A.7] \textbf{NO.} \hfill[1 point]

  From the previous result,
  \[ R_{max} = Re^{\frac{v^2}{4G\mu}}, \]
  we see that it is \textbf{not possible} to escape the gravitational
  field since for any speed, there is always a maximum distance
  $R_{max} < \infty$
\end{description}

\textbf{PART B:}

\begin{description}
\item[B.1] The energy density is given by
  \[ \rho = aT^4 \]
  \hfill[1 point]

  Equivalently:
  \[ \rho = \frac{4\sigma}{c}T^4
     = \frac{\pi^2 k_B^4}{15\,\hbar^3 c^3}T^4 \]
  \hfill[1 point]

  Note 3: This result arrives from the integration of the spectral energy
  density given by the Planck law over all frequencies. However, it is
  not required to show this integral.

\item[B.2]
  \[ r_0 = \frac{\hbar^{n_1}c^{n_2}}{k_B T} \]
  \begin{itemize}
  \item $r_0$ has dimensions of length: $[l]$
  \item $\hbar$ has dimensions of energy $\times$ time: $[e\,t]$
  \item $c$ has dimensions of speed: $[l/t]$
  \item $k_B T$ has dimensions of energy: $[e]$
  \end{itemize}
  Then:
  \[ [l] = \frac{[e\,t]^{n_1}\left[\frac{l}{t}\right]^{n_2}}{[e]} \]

  Since the LHS and the RHS should have the same dimensions, we get:
  \begin{align*}
  [l]&:\ 1 = n_2\\
  [e]&:\ 0 = n_1 - 1\\
  [t]&:\ n_1 = n_2
  \end{align*}
  \hfill[2 points]

  The solution is
  \[ n_1 = 1 \]
  \hfill[1 point]
  \[ n_2 = 1 \]
  \hfill[1 point]

\item[B.3] The energy of a piece of the string of length $L$ is
  \[ E = \rho\pi r_0^2 L = Mc^2 \]
  \hfill[1 point]

  where $M = \mu L$ is the mass of the piece. Solving for $\mu$:
  \[ \mu = \frac{\rho\pi r_0^2}{c^2} \]
  \hfill[1 point]

\item[B.4] The weak field condition is
  \[ \frac{2G\mu}{c^2}\ll 1 \]
  \[ 2G\frac{\rho\pi r_0^2}{c^2}\ll 1 \]
  \hfill[1 point]
  \[ 2G\frac{\pi}{c^2}\left(aT^4\right)
     \left(\frac{\hbar c}{k_B T}\right)^2\ll 1 \]
  \hfill[1 point]
  \[ \frac{2Ga\hbar^2\pi}{c^2 k_B^2}\,T^2\ll 1 \]
  \hfill[2 point]

  Where we used: $\mu = \frac{\rho\pi r_0^2}{c^2}$, $\rho = aT^4$,
  $r_0 = \frac{\hbar c}{k_B T}$
  \begin{align*}
  \frac{2Ga\hbar^2\pi}{c^2 k_B^2}\,T^2 &\ll 1\\
  \frac{2G\hbar^2\pi}{c^2 k_B^2}
    \left(\frac{\pi^2 k_B^4}{15\,\hbar^3 c^3}\right)T^2 &\ll 1\\
  \frac{2\pi^3}{15}\left(\frac{Gk_B^2}{\hbar c^5}\right)T^2 &\ll 1\\
  \frac{2\pi^3}{15}\,\frac{T^2}{T_{Pl}^2} &\ll 1
  \end{align*}
  \hfill[2 point]

  where $T_{Pl} = 1.416784\times 10^{32}\,K$ (known as the Planck
  Temperature)

  Note 4: the numerical factor $\frac{2\pi^3}{15}\sim 4.13$,
  \emph{(a 5\% error tolerance is accepted)}

  Note 5: From \textbf{A.4} we get that
  \[ 2G\mu = v^2 \]
  where $v$ is the speed with which a particle would orbit a string. Thus
  the weak field is rewritten as
  \[ \frac{v^2}{c^2}\ll 1 \]
  which is equivalent to a non-relativistic condition.

\item[B.5] The weak field condition is
  \[ 4.13\left(\frac{T}{T_{Pl}}\right)^2\ll 1 \]
  equivalently
  \[ \sim 2\frac{T}{T_{Pl}}\ll 1 \]
  \hfill[1 point]
  \begin{description}
  \item[i.] $\frac{2T_{EW}}{T_{Pl}}\sim
    2\times\frac{10^{15}}{10^{32}}\approx 1.4\times 10^{-17}\ll 1$
    \hfill[1 point]
  \item[ii.] $\frac{2T_{GUT}}{T_{Pl}}\sim
    2\times\frac{10^{29}}{10^{32}}\approx 1.4\times 10^{-3}\ll 1$
    \hfill[1 point]
  \end{description}

  Note 6: Using the following expression
  $4.13\left(\frac{T}{T_{Pl}}\right)^2\ll 1$

  One gets
  \begin{description}
  \item[i.] $4.13\left(\frac{T_{EW}}{T_{Pl}}\right)^2
    \sim 2.1\times 10^{-34}\ll 1$
  \item[ii.] $4.13\left(\frac{T_{GUT}}{T_{Pl}}\right)^2
    \sim 2.1\times 10^{-6}\ll 1$
  \end{description}
  And are also acceptable answers \emph{(a 5\% error tolerance is
  accepted)}

\item[B.6]\
  \begin{description}
  \item[i.] Yes \hfill[0.5 point]
  \item[ii.] Yes \hfill[0.5 point]
  \end{description}
\end{description}

\textbf{PART C:}

\begin{description}
\item[C.1] From the figure, the asymptotic condition for the observer to
  see a second image is that the light ray travelling from O directly to
  S should bend and travel along SE. This is the maximum angle of
  bending.

  As all angles are small, we can safely use
  \[ \sin x \approx \tan x \approx x \]
  Now,
  \[ \frac{p}{SP} < \delta\phi \]
  \hfill[2 points]
  \begin{align*}
  &< \frac{4\pi G\mu}{c^2}\\
  &< 2\pi\left(\frac{2G\mu}{c^2}\right)\\
  &< 2\pi\left(\frac{2\pi^3}{15}\right)
     \left(\frac{T^2}{T_{Pl}^2}\right)
  \end{align*}
  \hfill[2 points]

  And $SP \approx \left(D_{OE} - D_{ES}\right)$
  \[ p < \frac{4\pi^4}{15T_{Pl}^2}\,T^2\left(D_{OE} - D_{ES}\right) \]
  \hfill[2 points]

  \ioaafig{2021_TH_Q15-s3.pdf}

\item[C.2] In the figure, the blue lines represent the bending of two
  light rays corresponding to the images O1 and O2. Note that
  $D_{ES} \approx D_{ES1} \approx D_{ES2}$ and
  $D_{OS} \approx D_{OS1} \approx D_{OS2}$. Thus the angle
  $S1EO \approx S2EO \equiv \alpha$ and
  $S1OE \approx S2OE \equiv \beta$. $2\alpha$ is the angular separation
  we are looking for. \hfill[2 points]

  Further, notice
  \[ \alpha + \beta = \delta\phi \]
  \hfill[1 point]

  Using law of sines in the triangle \textbf{EOS1}
  \[ \frac{\alpha}{D_{S1O}} = \frac{\beta}{D_{ES1}} \]
  \hfill[1 point]

  we get
  \[ 2\alpha = 2\delta\phi
     \left(\frac{D_{OE} - D_{ES}}{D_{OE}}\right) \]
  \hfill[2 points]

\item[C.3] If $D_{OE} = 2D_{ES}$,
  \[ 2\alpha = \delta\phi
     = 2\pi\left(\frac{2\pi^3}{15}\right)
       \left(\frac{T_{GUT}^2}{T_{Pl}^2}\right)
     \approx 1.29\times 10^{-5} \]
  \hfill[2 points]

  and
  \[ \delta\phi = 1.22\frac{\lambda}{D} \]

  Thus, using
  $\lambda\in[3\times 10^{-7}m,\ 8\times 10^{-7}m]$
  \[ D = 1.22\frac{\lambda}{\delta\phi}
     \in\left[3.75\times 10^{-2}m,\ 7.51\times 10^{-2}m\right] \]
  \hfill[2 points]

  Any answer within this interval is valid.

  A remarkable result of this model is light deflection by a cosmic
  string, which leads to the possibility of detection through
  gravitational lensing. For instance, cosmic string trings moving across
  the line of sight with respect to an Earth-based observer will cause
  line-like discontinuities as shown in the figure below (taken from
  [1]).
  % sic: "cosmic string trings" as in source
  \ioaafig{2021_TH_Q15-s4.pdf}

  [1] M. V. Sazhin, O. S. Khovanskaya, M. Capaccioli, G. Longo,
  M. Paolillo, G. Covone, N. A. Grogin, E. J. Schreier, Gravitational
  lensing by cosmic strings: what we learn from the CSL-1 case,
  \emph{Monthly Notices of the Royal Astronomical Society}, Volume 376,
  Issue 4, March 2007, Pages 1731--1739.
\end{description}

\solhead{13.20}{Journey Between Galaxies}
% source id: 2022_TH_4
\begin{description}
\item[(a)] By the Hubble's Law, the rate of recession is
  \[ v(t) = \frac{\Delta d}{\Delta t} = H d(t) \]
  \hfill(1.0)

  Thus, by referring to the given relation,
  \[ d(t) = C e^{Ht} \]
  at time $t = 0$, $d(0) = d_0 = C$. Thus, the distance at a time $t$:
  \[ d(t) = d_0 e^{Ht} \]
  \hfill(1.0)

\item[(b)] It is easier to solve this problem, if we consider coordinate
  system that changes its scale in such a way that distance between earth
  and our destination does not change (called co-moving frame). Let $v(t)$
  be the velocity in this frame at time $t$ and let $l(t)$ be the total
  distance traveled until the moment $t$.
  \begin{align*}
  v(t)d(t) &= v_0 d_0\\
  v(t) &= v_0\frac{d_0}{d(t)} = v_0\frac{d_0}{d_0 e^{Ht}}\\
  v(t) &= \frac{\Delta l}{\Delta t} = -v_0 e^{-Ht}
  \end{align*}
  \hfill(3.0)
  \[ l(t) = -\frac{v_0}{-H}e^{-Ht} + C \]
  for specific cases:
  \[ l(t_0) = \frac{v_0}{H}e^{-Ht_0} + C \]
  \[ l(0) = \frac{v_0}{H} + C \]
  subtracting yields:
  \[ l(0) = \frac{v_0}{H} - \frac{v_0}{H}e^{-Ht_0} \]
  % sic: source writes l(0) for the difference l(0) - l(t_0)
  thus:
  \[ \therefore\ l_{\mathrm{Total}} = d_0 = -\frac{v_0}{H}(e^{-Ht_0} - 1) \]
  \hfill(2.0)

  Rearranging,
  \[ t_0 = -\frac{1}{H}\ln\left(1 - \frac{H d_0}{v_0}\right) \]
  \hfill(1.0)

  Thus, condition for reaching the planet at all:
  \begin{align*}
  1 - \frac{H d_0}{v_0} &> 0\\
  v_0 &> H d_0 = 70\ \mathrm{km/s}
  \end{align*}
  \hfill(1.0)

  (However at this speed it will take eternity to reach it)

  For $v_0 = 1000$\,km/s we can reach the planet, in:
  \[ t_0 \approx 5.3\times 10^{16}\,\mathrm{s} = 17\,\mathrm{Gyr} \]
  \hfill(1.0)
\end{description}

\solhead{13.21}{Dark Energy}
% source id: 2023_TH_Q5
Substituting the values of $H_0$ and $G$ from the table of constants,
\[ \rho_{\mathrm{cr}} = 9.202\times 10^{-27}\,\mathrm{kg\,m^{-3}}. \]
\hfill(1 point)

In flat geometry we have:
\[ \rho_{\mathrm{m},0} + \frac{E_\Lambda}{c^2} = \rho_{\mathrm{cr}}. \]
\hfill(2 points)

Hence $E_\Lambda$ is given by:
\[ E_\Lambda = \left[\rho_{\mathrm{cr}} - \rho_{\mathrm{m},0}\right]c^2
   \approx 5.756\times 10^{-10}\,\mathrm{J\,m^{-3}}. \]
\hfill(1 points)

The linear scale of the Universe, $a$, is related to the cosmological
redshift:
\[ a(z) = a_0/(1 + z). \]
\hfill(2 points)

Thus, the matter density, $\rho_{\mathrm{m}}$, increases with redshift:
\[ \rho_{\mathrm{m}}(z) = \rho_{\mathrm{m},0}\,(1 + z)^3. \]
\hfill(2 points)

We substitute the matter density, $\rho_{\mathrm{m}}$ by the energy
density, $E_{\mathrm{m}}$, using the relationship $E = mc^2$
\[ E_{\mathrm{m}}(z) = \rho_{\mathrm{m}}(z)\,c^2
   = \rho_{\mathrm{m},0}\,(1 + z)^3\,c^2. \]
\hfill(2 points)

We finally get
\[ z_{\mathrm{eq}} = \left(\frac{E_\Lambda}{\rho_{\mathrm{m},0}\,c^2}\right)^{1/3} - 1
   \approx 0.32. \]
\hfill(2 points)

\solhead{13.22}{CMB}
% source id: 2024_TH_T5
\begin{description}
\item[(a)] In order to find the redshift, we need $\lambda_{\mathrm{today}}$.
  Since $T_{\mathrm{today}}$ refers to the temperature at the emission
  peak, we use Wien's law:
  \begin{align*}
  \lambda_{\mathrm{today}}\cdot T_{\mathrm{today}} &= b\\
  \lambda_{\mathrm{today}} &\approx 1\,\mathrm{mm}
  \end{align*}
  \hfill(1.0)

  So the redshift of the CMB, when it is emmited at the infrared spectrum
  (with the values mentioned) is
  \[ z = \frac{\lambda_{\mathrm{today}} - \lambda_{\mathrm{infrared}}}
              {\lambda_{\mathrm{infrared}}} \]
  \hfill(1.0)
  \[ z \approx \frac{1 - 0.1}{0.1} = 9 \]
  \hfill(1.0)

  \textbf{Note:} It is also possible to find $T_{\mathrm{infrared}}$ using
  Wien's law, and then use the definition of redshift to find $z$, which
  leads to the same answer.

\item[(b)] As the Universe is assumed as composed of just matter (i.e.
  $\Omega_m = 1$), so the expression for the age of the Universe is given,
  in terms of the scale factor $\left(a = \frac{1}{1+z}\right)$, by
  \[ a_{\mathrm{ir}} = \frac{1}{1+z} = \frac{1}{1+9} = 0.1 \]
  \hfill(1.0)
  \[ \text{Now, }\ \frac{\mathrm{d}a}{\mathrm{d}t} = H(t)\cdot a(t) \]
  \hfill(1.0)

  And, from the first Friedmann equation for a matter-dominated universe,
  we know that:
  \[ H^2 = \frac{H_0^2}{a^3} \]
  \hfill(2.0)

  Therefore:
  \begin{align*}
  \frac{da}{dt} &= \frac{H_0}{a^{1.5}(t)}a(t) = \frac{H_0}{\sqrt{a(t)}}\\
  \therefore\ t &= \int_0^{0.1}\frac{\sqrt{a}}{H_0}\,\mathrm{d}a\\
  &= \frac{1}{H_0}\int_0^{0.1}\sqrt{a}\ \mathrm{d}a
   = \frac{2}{3\,H_0}\,a^{1.5}\big|_0^{0.1}
  \end{align*}
  \hfill(2.0)
  \begin{align*}
  &= \frac{2\times 0.1^{1.5}}{3\times 70}\\
  &= 3\times 10^{-4}\,\mathrm{s\,Mpc/km}\\
  &= \frac{3\times 10^{-4}\times 3.086\times 10^{22}}
          {1000\times 86400\times 365.2422}
  \end{align*}
  \[ t \approx 0.3\,\mathrm{Gyr} \]
  \hfill(1.0)
\end{description}

\solhead{13.23}{Big Bang Nucleosynthesis}
% source id: 2025_TH_T10
\begin{description}
\item[(T10.1)] The ratio of the number density of protons to that of
  neutrons at $t_{\mathrm{wk}}$ will be given by
  \[ \frac{N_{\mathrm{n}}(t_{\mathrm{wk}})}{N_{\mathrm{p}}(t_{\mathrm{wk}})}
     = \left(\frac{m_{\mathrm{n}}}{m_{\mathrm{p}}}\right)^{3/2}
       \exp\left(-\frac{(m_{\mathrm{n}} - m_{\mathrm{p}})c^2}
                        {k_{\mathrm{B}}T_{\mathrm{wk}}}\right) \]
  \hfill(2.0)

  Substituting the values of the masses of neutron and proton, and
  $k_{\mathrm{B}}T_{\mathrm{wk}} = 800$\,keV, one obtains
  \[ X_{\mathrm{n}}(t_{\mathrm{wk}})
     = \frac{N_{\mathrm{n}}}{N_{\mathrm{p}} + N_{\mathrm{n}}} = 0.166 \]
  \hfill(2.0)

\item[(T10.2a)] For a black body spectrum,
  $\lambda_{\mathrm{max}}T$ should be equal to a constant. \hfill(1.0)

  Since the wavelength of light is $\lambda \propto a$, this implies
  $T \propto a^{-1}$, which implies $\beta = -1$. \hfill(1.0)

  Alternatively, for black body radiation, energy density is proportional
  to $T^4$, but the energy density is proportional to $a^{-4}$.
  Therefore, $T \propto a^{-1}$.

\item[(T10.2b)] Graph B should be ticked. Wien's law removes options D
  and E, Planck curves at different temperatures should not intersect,
  and so C and A are also incorrect.

\item[(T10.3a)] Each of the 4 points correctly located on the graph
  \hfill(1.0)

  Drawing a smooth curve passing through all the points (within one grid
  point) \hfill(1.0)

  Non-smooth lines, or straight lines not passing through all points get
  0 credit.
  \ioaafig{2025_TH_T10-s1.pdf}

\item[(T10.3b)] The value of
  $k_{\mathrm{B}}T_{\mathrm{nuc}} = 66.0\,\mathrm{keV}$. \hfill(1.0)

  Full credit range: 65.4\,keV to 66.6\,keV.

\item[(T10.3c)] As $T \propto a^{-1}$, and $a \propto t^{1/2}$, we have
  $T \propto t^{-1/2}$. \hfill(2.0)

  The temperature and time when all neutrons and protons fuse into
  Helium are $T_{\mathrm{nuc}}$ (66.0\,keV) and $t_{\mathrm{nuc}}$,
  respectively.
  \[ t_{\mathrm{nuc}} = t_{\mathrm{wk}}
     \left(\frac{T_{\mathrm{wk}}}{T_{\mathrm{nuc}}}\right)^2
     = 250\,\mathrm{s} \]
  \hfill(2.0)

\item[(T10.4)] During $t_{\mathrm{wk}} < t < t_{\mathrm{nuc}}$, as the
  Universe cools from 800\,keV to 66.0\,keV, neutrons decay into protons
  and with a half life of 610.4\,s, which corresponds to a mean life time
  $\tau = 880.6$\,s. \hfill(1.0)

  The quantity $(N_{\mathrm{p}} + N_{\mathrm{n}})a^3$ does not change as
  the increase in protons is compensated by the decrease in neutrons.
  \hfill(1.0)

  However, $N_{\mathrm{n}}a^3$ will go down by an exponential factor
  compared to its value computed in the first question as the neutrons
  decay into protons, i.e.,
  \[ X_{\mathrm{n,\,nuc}} = X_{\mathrm{n,\,wk}}
     \exp(-(t_{\mathrm{nuc}} - t_{\mathrm{wk}}/880.6)) \]
  % sic: misplaced parenthesis in source; the numerical line below has
  % the intended (t_nuc - t_wk)/880.6
  \hfill(2.0)
  \[ X_{\mathrm{n,\,nuc}} = X_{\mathrm{n,\,wk}}
     \exp(-(250 - 1.70)/880.6) = 0.125 \]
  \hfill(1.0)

  Give part credit of 1 marks if half life time is used instead of the
  mean life time in the denominator of the exponential.

\item[(T10.5)] Helium forms with the combination of 2 protons and 2
  neutrons (ignoring mass deficit), and thus a number density
  $N_{\mathrm{n}}$ of neutrons and correspondingly same number density of
  protons will be locked in Helium. The number density of both species
  goes as $a^{-3}$, and so will cancel from the numerator and
  denominator. \hfill(1.0)

  Even after Helium is formed the sum of the mass of Helium and Hydrogen
  will still be the mass of all the neutrons and protons. \hfill(1.0)

  Thus, the abundance of Helium by mass should be
  \[ Y_{\mathrm{prim}}
     = \frac{(m_{\mathrm{p}} + m_{\mathrm{n}})N_{\mathrm{n}}}
            {m_{\mathrm{p}}N_{\mathrm{p}} + m_{\mathrm{n}}N_{\mathrm{n}}}
     \approx 2X_{\mathrm{n,\,nuc}} = 0.250 \]
  \hfill(1.0)

  Alternatively, the value of $X_{\mathrm{n,\,nuc}}$ is $1/8$. After
  formation of Helium all neutrons will combine with protons. Therefore,
  the mass in Helium will be $2/8$ of the total mass.

\item[(T10.6a)] By fitting a line to these data points the Helium
  abundance for the said galaxy would be: $Y = 0.2480$ \hfill(2.0)

  Full credit range: 0.2475 to 0.2485.
  \ioaafig{2025_TH_T10-s2.pdf}

\item[(T10.6b)] The slope of the line yields
  $dY/d(\mathrm{O/H}) = 23$. \hfill(2.0)

  Full credit range: 19 to 27.

\item[(T10.6c)] The primordial Helium abundance,
  $Y_{\mathrm{prim}}^{\mathrm{obs}} = 0.2440$. \hfill(2.0)

  Full credit range: 0.2435 to 0.2445.

\item[(T10.7)] From the Saha equation, we see that
  $N_{\mathrm{D}}/N_{\mathrm{n}}(T) \propto \eta$, thus when
  $\eta\downarrow$, $N_{\mathrm{D}}/N_{\mathrm{n}}(T)$ will also
  $\downarrow$. From the plot in (T10.3a), this would mean that
  $N_{\mathrm{D}}/N_{\mathrm{n}}(T)$ would be unity at a lower value of
  $T_{\mathrm{nuc}}$ which will go $\downarrow$. A lower temperature
  corresponds to a longer $t_{\mathrm{nuc}}$ ($\uparrow$). This would
  result in more neutrons decaying into protons and thus less
  $Y_{\mathrm{prim}}$ $\downarrow$.

  \begin{center}
  \begin{tabular}{ll}
  $\eta$ & $\downarrow$\\
  $N_{\mathrm{D}}/N_{\mathrm{n}}(T)$ & $\downarrow$\\
  $T_{\mathrm{nuc}}$ ($N_{\mathrm{D}} = N_{\mathrm{n}}$) & $\downarrow$\\
  $t_{\mathrm{nuc}}$ & $\uparrow$\\
  $X_{\mathrm{n,\,nuc}}$ & $\downarrow$\\
  $Y_{\mathrm{prim}}$ & $\downarrow$
  \end{tabular}
  \end{center}

  Each correct box fetches 0.5 marks. We give a bonus 0.5 marks if all 5
  boxes are marked correctly.
\end{description}

\solhead{13.24}{Virgo Cluster}
% source id: 2008_DA_D1
\begin{description}
\item[1.] From table above, we have
  $\sum_i \dfrac{d_i}{s_i^2} = 56.1082$ and
  $\sum_i \dfrac{1}{s_i^2} = 3.5485$, so:
  \[ d_{avg} = \frac{56.1082}{3.5485}\ \mathrm{Mpc} = 15.81\ \mathrm{Mpc} \]
  \hfill(30\%)

\item[2.] The uncertainty (rms):
  \[ s_d = \sqrt{\frac{1}{\sum_i \frac{1}{s_i^2}}}
     = \sqrt{\frac{1}{3.5485}}\ \mathrm{Mpc} = 0.53\ \mathrm{Mpc} \]
  \[ d_{\mathrm{Virgo}} = 15.81 \pm 0.53\ \mathrm{Mpc} \]
  \hfill(10\%)

\item[3.] Hubble constant
  \[ H_0 = \frac{v}{d_{avg}} = \frac{1136\ \mathrm{km/s}}{15.81\ \mathrm{Mpc}}
     = 71.9\ \mathrm{km/sec/Mpc} \]
  \hfill(15\%)

  Uncertainty:
  \[ s_{H_0} = \frac{s_d}{d_{avg}}H_0 = \frac{0.53}{15.81}\times 71.85
     = 0.0335\times 71.85 = 2.4\ \mathrm{km/sec/Mpc} \]
  \hfill(15\%)

  \textbf{Hubble constant:} $H_0 = 72 \pm 2\,km/\sec/Mpc$
  \hfill(5\%)

\item[4.] Hubble time:
  \[ H_0^{-1} = \frac{1}{H_0} = \frac{1}{71.85\ \mathrm{km/sec/Mpc}}
     \times 3.086\times 10^{10}\ \mathrm{km/Mpc}
     = 4.29\times 10^{17}\,\mathrm{sec} = 13.6\,\mathrm{Gyr} \]
  \hfill(20\%)

  Uncertainty is 3.35\% which is 0.45 Gyr.
  \hfill(5\%)
\end{description}

\solhead{13.25}{The search for dark matter}
% source id: 2017_DA_D2
\begin{description}
\item[(D2.1)] Redshift:
  $z = \dfrac{\lambda_{\mathrm{obs}} - \lambda_{\mathrm{emit}}}
             {\lambda_{\mathrm{emit}}}$ \hfill[0.5]

  Observed rotation velocity: $V_{\mathrm{obs}} = zc$ \hfill[0.5]

  \begin{center}
  \begin{tabular}{cccccc}
  $r$ (kpc) & $z$ & $V_{\mathrm{obs}}$ (km/s) & $V_{\mathrm{DM}}$ (km/s)
    & $M(r)$ ($M_\odot$) & $M_{\mathrm{DM}}(r)$ ($M_\odot$)\\
  0.70 & $1.37\times 10^{-4}$ & 41.11 & 35.25
    & $2.75\times 10^{8}$ & $2.02\times 10^{8}$\\
  1.40 & $2.29\times 10^{-4}$ & 68.52 & 60.10
    & $1.53\times 10^{9}$ & $1.18\times 10^{9}$\\
  2.09 & $2.79\times 10^{-4}$ & 83.60 & 71.52
    & $3.40\times 10^{9}$ & $2.49\times 10^{9}$\\
  2.79 & $2.96\times 10^{-4}$ & 88.62 & 72.53
    & $5.09\times 10^{9}$ & $3.41\times 10^{9}$\\
  3.49 & $3.00\times 10^{-4}$ & 89.99 & 70.77
    & $6.57\times 10^{9}$ & $4.06\times 10^{9}$\\
  4.89 & $3.09\times 10^{-4}$ & 92.73 & 74.25
    & $9.78\times 10^{9}$ & $6.27\times 10^{9}$\\
  6.25 & $3.05\times 10^{-4}$ & 91.36 & 73.65
    & $1.21\times 10^{10}$ & $7.88\times 10^{9}$\\
  7.10 & $3.05\times 10^{-4}$ & 91.36 & 73.03
    & $1.38\times 10^{10}$ & $8.80\times 10^{9}$\\
  9.03 & $3.06\times 10^{-4}$ & 91.82 & 75.54
    & $1.77\times 10^{10}$ & $1.20\times 10^{10}$\\
  12.05 & $3.06\times 10^{-4}$ & 91.82 & 78.74
    & $2.36\times 10^{10}$ & $1.74\times 10^{10}$
  \end{tabular}
  \end{center}

\item[(D2.2)]
  \ioaafig{2017_DA_D2-s1.pdf}
  From the official graph, the curves in descending order of contribution
  to the observed velocity are $V_{\mathrm{obs}}$, $V_{\mathrm{DM}}$,
  $V_{*}$, $V_{\mathrm{gas}}$.

\item[(D2.3)] Evaluating for $\rho_0$:
  \[ M_{\mathrm{DM}}(r) = 4\pi\rho_0 r_C^2
     \left[r - r_C\arctan(r/r_C)\right] \]
  \[ M_{DM}(r) = 4\pi\rho_0 r_C^3\left[x - \arctan(x)\right],
     \text{ where } x = r/r_C \]
  \hfill[1.0]
  \[ M_{DM}(r) \approx 4\pi\rho_0 r_C^3
     \left[x - \left(x - \frac{x^3}{3}\right)\right],
     \text{ for small } x \]
  \hfill[1.0]
  \[ M_{DM}(r) \approx 4\pi\rho_0 r_C^3\left(\frac{x^3}{3}\right)
     = \frac{4\pi\rho_0 r^3}{3} \]
  \[ \rho_0 \approx \frac{2.02\times 10^{8}M_\odot\times 3}
     {4\pi(0.7\,kpc)^3} = 1.42\times 10^{8}\,M_\odot/kpc^3 \]
  \hfill[1.5]

  Selected correct data to put into the formula gets 0.5, correct answer
  gets 0.5 and correct unit gets 0.5.

  \[ M_{DM}(r) = 4\pi\rho_0 r_C^2
     \left[r - r_C\arctan\left(\frac{r}{r_C}\right)\right] \]

  Evaluating for $r_C$ (Method 1):
  \[ M_{DM}(r) \approx 4\pi\rho_0 r_C^2
     \left[r - r_C\frac{\pi}{2}\right], \]
  \hfill[1.0]

  Take the last two data points at large $r$, then we get (for
  $r \gg r_C$
  \[ \Delta M_{DM}(r) \approx 4\pi\rho_0 r_C^2\left[\Delta r\right] \]
  \hfill[1.5]
  \[ r_C = 1.01\ kpc \]
  \hfill[1.0]

  Evaluating for $r_C$ (Method 2):
  \[ M_{DM}(r) \approx 4\pi\rho_0 r_C^2
     \left[r - r_C\frac{\pi}{2}\right], \]
  \hfill[1.0]

  Take the last data point, we get a cubic equation which students can
  solve to give
  \[ r_C = -0.855,\ 0.964,\ 7.56\ \text{kpc} \]
  \hfill[1.0]

  Select the correct $r_C = 0.964$ for the final answer because for the
  last data point $r \gg r_C$ \hfill[1.5]

  Less accurate $r_C$:

  Omit the last term,
  \[ M_{DM}(r) \approx 4\pi\rho_0 r_C^2 r \]
  \hfill[1.0]
  \[ r_C \approx \sqrt{\frac{1.74\times 10^{10}M_\odot}
     {4\pi\times 1.40\times\frac{10^{8}M_\odot}{kpc^3}\times 12.05\,kpc}}
     = 0.901\ kpc \]
  \hfill[1.0]

\item[(D2.4)] Solution 1:
  \[ x = \left[r - r_c\arctan\left(\frac{r}{r_C}\right)\right] \]
  \[ y = M_{DM}(r) \]

  \begin{center}
  \begin{tabular}{cc}
  $x = \left[r - r_C\arctan(r/r_C)\right]$ (kpc)
    & $y = M_{\mathrm{DM}}(r)$ (kg)\\
  % sic: source table header says (kg); the values are in solar masses
  0.010855 & $2.02\times 10^{8}$\\
  0.0803 & $1.18\times 10^{9}$\\
  0.2386 & $2.49\times 10^{9}$\\
  0.495 & $3.41\times 10^{9}$\\
  0.838 & $4.06\times 10^{9}$\\
  1.718 & $6.27\times 10^{9}$\\
  2.738 & $7.88\times 10^{9}$\\
  3.428 & $8.80\times 10^{9}$\\
  5.093 & $1.20\times 10^{10}$\\
  7.854 & $1.74\times 10^{10}$
  \end{tabular}
  \end{center}

  \ioaafig[0.65]{2017_DA_D2-s2.pdf}
  From the linear fit,
  \[ \text{slope} = \frac{\Delta y}{\Delta x} \approx
     2\times 10^{9}\,M_\odot/kpc,\qquad
     \rho_0 = \frac{\text{slope}}{4\pi r_c^2}
     = 1.43\times 10^{8}\ M_\odot/kpc^3 \]

  Solution 2:
  \[ x = \frac{r_C}{r}\arctan\left(\frac{r}{r_C}\right) \]
  \[ y = M_{DM}(r) \]
  % sic: source labels y as M_DM(r); the tabulated values are M_DM(r)/r

  \begin{center}
  \begin{tabular}{cc}
  $x = \dfrac{r_C}{r}\arctan\left(\dfrac{r}{r_C}\right)$
    & $y = M_{\mathrm{DM}}(r)$ ($M_\odot$)\\
  0.984 & $2.89\times 10^{8}$\\
  0.943 & $8.40\times 10^{8}$\\
  0.886 & $1.19\times 10^{9}$\\
  0.822 & $1.22\times 10^{9}$\\
  0.760 & $1.16\times 10^{9}$\\
  0.649 & $1.28\times 10^{9}$\\
  0.562 & $1.26\times 10^{9}$\\
  0.517 & $1.24\times 10^{9}$\\
  0.436 & $1.33\times 10^{9}$\\
  0.348 & $1.44\times 10^{9}$
  \end{tabular}
  \end{center}

  \ioaafig[0.70]{2017_DA_D2-s3.pdf}

\item[(D2.5)] Solution 1: $r_c = 0.906\ kpc$ from D2.3

  \begin{center}
  \begin{tabular}{cc}
  $r$ (kpc) & $\ln\left[\rho_{\mathrm{DM}}(r)\right]$
    $\left(\frac{M_\odot}{kpc^3}\right)$\\
  0.70 & 16.11\\
  1.40 & 15.36\\
  2.09 & 14.73\\
  2.79 & 14.23\\
  3.49 & 13.81\\
  4.89 & 13.17\\
  6.25 & 12.69\\
  7.10 & 12.44\\
  9.03 & 11.97\\
  12.05 & 11.40
  \end{tabular}
  \end{center}
  \begin{center}
  \textbf{Graph Number D2.5: Solution 1}\\[0.25em]
  \includegraphics[width=0.57\linewidth]{2017_DA_D2-s4.pdf}
  \end{center}

  Solution 2: $r_c = 3.2\ kpc$

  \begin{center}
  \begin{tabular}{cc}
  $r$ (kpc) & $\ln\left[\rho_{\mathrm{DM}}(r)\right]$
    $\left(\frac{M_\odot}{kpc^3}\right)$\\
  0.70 & 12.79\\
  1.40 & 11.95\\
  2.09 & 11.47\\
  2.79 & 11.21\\
  3.49 & 11.07\\
  4.89 & 11.16\\
  6.25 & 11.26\\
  7.10 & 11.16\\
  9.03 & 11.27\\
  12.05 & 11.41
  \end{tabular}
  \end{center}
  \ioaafig[0.55]{2017_DA_D2-s5.pdf}
\end{description}

\solhead{13.26}{Scaling Relations}
% source id: 2021_DA_DA1
\textbf{Part 1.}

\begin{description}
\item[1.1)] From the plot on the right we can get the equation of the
  straight line as
  \[ y = -0.76\cdot x + 7.32 \]
  \hfill[3 points]

  which taking to the power of 10 is equivalent to:
  $f_g = 10^{7.32}M_*^{-0.76}$, being $f_g$ the gas fraction. Multiplying
  by the stellar mass we finally get:
  \[ M_{gas} = 10^{7.32}\times M_*^{0.24} \]
  \[ a = 10^{7.32} = 2.09\times 10^{7}\ \text{[1 point]};\qquad
     b = 0.24\ \text{[1 point]} \]

\item[1.2)] \hfill[15 points]

  \begin{center}
  \begin{tabular}{cccccc}
  V & LogV & K & Zx & Zy & ZxZy\\
  79.4 & 1.8998 & $-16.8$ & $-1.1778$ & 1.3873 & $-1.6339$\\
  100.1 & 2.0004 & $-19.2$ & $-0.7831$ & 0.4970 & $-0.3892$\\
  158.5 & 2.2000 & $-21.3$ & $-0.0001$ & $-0.2819$ & 0.0000\\
  251.2 & 2.4000 & $-21.4$ & 0.7845 & $-0.3190$ & $-0.2502$\\
  316.2 & 2.5000 & $-24$ & 1.1765 & $-1.2834$ & $-1.5100$\\
  \hline
  Ave & 2.2001 & $-20.54$ & & sum ZxZy & $-3.7833$\\
  St. Dev & 0.2549 & 2.696 & & r & $-0.9458$
  \end{tabular}
  \end{center}

  The best-fit straight line is
  $K = -10\cdot\log_{10}(V_{max}) + 1.47$.
\end{description}

\textbf{Part 2.}

\begin{description}
\item[2.1)] \hfill[10 points]

  Galaxies are at the same distance, therefore the difference in apparent
  magnitudes is the same as in absolute magnitudes, and this relates to a
  luminosity ratio:
  \[ (k_1 - k_2) = (K_1 - K_2)
     = -2.5\times\log_{10}\!\left(\frac{L_1}{L_2}\right) \]
  \hfill[2 points]

  For having the same stellar populations, and ignoring interstellar
  extinction, the luminosity of each galaxy is proportional to its
  stellar mass with the same mass-to-light ratio for both galaxies, such
  that:
  \[ (k_1 - k_2) = -2.5\times\log_{10}
     \!\left(\frac{M_{*1}}{M_{*2}}\right) \]
  \hfill[2 points]

  The difference in magnitudes being 6, we obtain:
  \[ \left(\frac{M_{*1}}{M_{*2}}\right) = 10^{2.4} \]
  \hfill[2 points]

\item[2.2)] \hfill[4 points]

  Using the last result in combination to that of 1A:
  \[ \frac{M_{gas1}}{M_{gas2}}
     = \left(\frac{M_{*1}}{M_{*2}}\right)^{0.24} \]
  \hfill[2 points]
  \[ \frac{M_{gas1}}{M_{gas2}} = 10^{0.576} \]
  \hfill[2 points]

\item[2.3)] \hfill[6 points]

  We found the slope of the TF relation in 1B:
  $\frac{\Delta K}{\Delta\log(V_{max})} = 10$, with $\Delta K = 6$ for
  our galaxies: $\frac{V_{max1}}{V_{max2}} = 10^{0.6}$ [2 points]. Now
  consider the mass distribution as approximately
  spherically-symmetric, giving:
  \[ V_{max} = \sqrt{\frac{G\cdot M_{tot}(<R_{max})}{R_{max}}}, \]
  \hfill[2 points]

  which combined with the ratio just found two lines above leads to:
  \[ \frac{M_{tot1}}{M_{tot2}} = 10^{1.2} \]
  \hfill[2 points]
\end{description}

\textbf{Part 3.} \hfill[15 points]

Baryonic mass is
$M_{Baryonic} = \frac{4.39\times 10^{11}}{7.82}\,[M_\odot]$ and Dark
matter mass is given by
$M_{dm} = 6.82\times\frac{4.39\times 10^{11}}{7.82}\,[M_\odot]$. The
baryonic contribution needs to be decomposed into stellar and gaseous
masses, which is not doable analytically but must be faced numerically,
by tuning the numbers, knowing their sum and knowing the relation between
them from 1.1. The final table looks like this:

\begin{center}
\begin{tabular}{cccccc}
Galaxy & Apparent magnitude $k$ & $M_{gas}$ [$M_\odot$]
  & $M_{*}$ [$M_\odot$] & $M_{dm}$ [$M_\odot$] & $M_{tot}$ [$M_\odot$]\\
$G_1$ & 19.2 & $7.67\times 10^{9}$ & $4.85\times 10^{10}$
  & $3.83\times 10^{11}$ & $4.39\times 10^{11}$
\end{tabular}
\end{center}

Equation to be solved: $1 - k = 0.1415\,k^{0.24}$. This gives
$k = 0.8634$

\textbf{Part 4}

\begin{description}
\item[4.1)] \hfill[4 points]

  We basically found in 2.1 that
  $\left(\frac{M_{*1}}{M_{*2}}\right) = 10^{\frac{\Delta K}{-2.5}}$,
  being $\Delta K = -6$. In this new case, we have the extreme values of
  $\Delta K$ to be $-6.4$ and $-5.6$, therefore the extreme values for
  the exponent are these values divided by $-2.5$, then
  $e\in[2.24,\ 2.56]$

\item[4.2)] \hfill[10 points]

  We have to invert the equation from 1.2 to get:
  $\log_{10}(V_{max}) = -\frac{K}{10} + 0.1467$. Now we can get the
  difference between the measured values and linear-fit predictions,
  approximating numbers to the second digit:

  \begin{center}
  \begin{tabular}{ccccc}
  $K$ [mag] & $V_{max}$ [km/s] & $\log_{10}(V_{max})$ measured
    & $\log_{10}(V_{max})$ from linear fit
    & $|\Delta\log_{10}(V_{max})|$\\
  $-16.8$ & 79.4 & 1.9 & 1.83 & 0.07\\
  $-19.2$ & 100.1 & 2.0 & 2.07 & 0.07\\
  $-21.3$ & 158.5 & 2.2 & 2.28 & 0.08\\
  $-21.4$ & 251.2 & 2.4 & 2.29 & 0.11\\
  $-24.0$ & 316.2 & 2.5 & 2.55 & 0.05
  \end{tabular}
  \end{center}

  Taking two times the RMS of the residuals:
  $\sigma_{stat} = 0.157$

\item[4.3)] With the uncertainties, the expression
  $\log_{10}(V_{max}) = -\frac{K}{10} + 0.1467$ becomes:
  \[ \log_{10}(V_{max}) = -\frac{K}{10} + 0.1467\pm\frac{0.2}{10}
     \pm 0.157 \]
  \hfill[2 points]

  Writing the same expression down for galaxy 1 and 2, subtracting, and
  using $\Delta K = -6$ we get:
  \[ \log_{10}\!\left(\frac{V_{max1}}{V_{max2}}\right)
     = 0.6\pm 0.02\pm 0.157\pm 0.02\pm 0.157 \]
  \hfill[4 points]

  Which leads to the extreme values: $0.954$ and $0.246$.
  \hfill[2 points]

  Finally remembering that circular velocities scale as the square root
  of the total mass, we have to multiply this exponent by 2, getting the
  final answer:
  \[ g\in[0.49,\ 1.91] \]
  \hfill[2 points]
\end{description}

\solhead{13.27}{Galactic Surveys}
% source id: 2022_DA_2
\textbf{Solutions -- Part 1}

\begin{description}
\item[(a)] If we plot given coordinates in XY, XZ and YZ planes we'll see
  the following pictures. See Figures 2, 3 and 4 (0.5p for each graph).
  \ioaafig{2022_DA_2-s1.pdf}
  \ioaafig{2022_DA_2-s2.pdf}
  \ioaafig{2022_DA_2-s3.pdf}
  By observing Figure 2 we can deduce that the group is edge-on and its
  normal lies in XY plane. $X'$ axis is along line of sight so it can be
  calculated by normalizing
  $8.3\hat{i} + 6.2\hat{j} + 1.1\hat{k}$. So
  \[ \hat{x}' = 0.79667866\,\hat{i} + 0.59510936\,\hat{j}
     + 0.10558392\,\hat{k} \]
  (0.5p) $Z'$ axis lays in XY plane and is perpendicular to $X'$ so it
  can be calculated by normalizing $\hat{x}'\times\hat{k}$
  \[ \hat{z}' = 0.59845448\,\hat{i} - 0.80115681\,\hat{j} \]
  \hfill(2p)

\item[(b)] One can choose $x'$ as line of sight (found in part (a)) (1p):
  \[ \hat{x}' = 0.79667866\,\hat{i} + 0.59510936\,\hat{j}
     + 0.10558392\,\hat{k} \]
  $z'$ can be chosen as halo disk normal (found in part (a)) (1p):
  \[ \hat{z}' = 0.59845448\,\hat{i} - 0.80115681\,\hat{j} \]
  And finally for $y'$ (1p):
  \[ \hat{y}' = \hat{x}'\times\hat{z}'
     = 0.08458928\,\hat{i} + 0.06318717\,\hat{j} - 0.9944104\,\hat{k} \]

\item[(c)] Each coordinate in $X'Y'Z'$ system can be found by following
  equations
  \begin{align*}
  x' &= \hat{x}'\cdot(\vec{r} - \vec{r}_c)\\
  y' &= \hat{y}'\cdot(\vec{r} - \vec{r}_c)\\
  z' &= \hat{z}'\cdot(\vec{r} - \vec{r}_c)
  \end{align*}
  Where $\hat{x}'$, $\hat{y}'$ and $\hat{z}'$ are unit vectors of local
  $X'Y'Z'$ system, $\vec{r}$ is coordinate of a galaxy and $\vec{r}_c$
  is coordinate of the center.

  This values are calculated in the following table (0.5p for each
  correct $x'$, $y'$, $z'$ tuple):

  \begin{center}
  \begin{tabular}{rrrrrrr}
  & $x - x_c$ & $y - y_c$ & $z - z_c$ & $x'$ & $y'$ & $z'$\\
  1 & 0.101124 & 0.109271 & 0.294476 & 0.176684 & $-0.277372$ & $-0.027025$\\
  2 & 1.582909 & 0.989392 & 0.405895 & 1.892722 & $-0.207212$ & 0.154641\\
  3 & 0.582735 & 0.213285 & 1.126426 & 0.710112 & $-1.057359$ & 0.177866\\
  4 & 0.144328 & 0.369381 & 1.338567 & 0.476137 & $-1.295537$ & $-0.209558$\\
  5 & $-0.517770$ & $-0.151687$ & $-1.458202$ & $-0.656729$ & 1.396668 & $-0.188336$\\
  6 & 0.480059 & 0.105110 & 1.362674 & 0.588881 & $-1.307808$ & 0.203084\\
  7 & $-0.309855$ & $-0.319116$ & $-0.568033$ & $-0.496739$ & 0.518483 & 0.070227\\
  8 & 0.239541 & 0.188218 & $-0.730736$ & 0.225694 & 0.758807 & $-0.007438$\\
  9 & $-0.324800$ & $-0.169551$ & 0.116378 & $-0.347375$ & $-0.153915$ & $-0.058541$\\
  10 & $-1.227659$ & $-0.768684$ & $-0.612379$ & $-1.500158$ & 0.456538 & $-0.118862$\\
  11 & $-0.263209$ & $-0.234869$ & $-0.485227$ & $-0.400698$ & 0.445410 & 0.030649\\
  12 & 0.381200 & 0.282515 & $-0.365581$ & 0.433222 & 0.413634 & 0.001792\\
  13 & $-1.185021$ & $-0.941303$ & 0.302593 & $-1.472310$ & $-0.460620$ & 0.044950\\
  14 & 1.286640 & 0.930147 & $-0.777889$ & 1.496445 & 0.941150 & 0.024802\\
  15 & $-0.106760$ & $-0.095804$ & $-0.184997$ & $-0.161600$ & 0.168879 & 0.012863\\
  16 & $-0.686865$ & $-0.352750$ & 0.566244 & $-0.697350$ & $-0.643470$ & $-0.128449$
  \end{tabular}
  \end{center}

\item[(d)] Spherical coordinates can by calculated by following
  relations:
  \[ r = \sqrt{x'^2 + y'^2 + z'^2} \]
  \[ \phi = \arctan\frac{y'}{x'}\ \text{and signs of $x'$ and $y'$} \]
  \[ \theta = \arccos\frac{z}{r} \]
  Calculated values are given in the following table (0.5p for each
  correct $r$, $\phi$, $\theta$ tuple):

  \begin{center}
  \begin{tabular}{rrrrrrr}
  & $x'$ & $y'$ & $z'$ & $r$ & $\phi$ & $\theta$\\
  1 & 0.176684 & $-0.277372$ & $-0.027025$ & 0.329974 & 5.279566 & 1.652789\\
  2 & 1.892722 & $-0.207212$ & 0.154641 & 1.910301 & 6.174141 & 1.489757\\
  3 & 0.710112 & $-1.057359$ & 0.177866 & 1.286042 & 5.303793 & 1.432047\\
  4 & 0.476137 & $-1.295537$ & $-0.209558$ & 1.396079 & 5.064586 & 1.721471\\
  5 & $-0.656729$ & 1.396668 & $-0.188336$ & 1.554814 & 2.010330 & 1.692226\\
  6 & 0.588881 & $-1.307808$ & 0.203084 & 1.448581 & 5.135477 & 1.430138\\
  7 & $-0.496739$ & 0.518483 & 0.070227 & 0.721461 & 2.334780 & 1.473302\\
  8 & 0.225694 & 0.758807 & $-0.007438$ & 0.791695 & 1.281697 & 1.580191\\
  9 & $-0.347375$ & $-0.153915$ & $-0.058541$ & 0.384430 & 3.558678 & 1.723672\\
  10 & $-1.500158$ & 0.456538 & $-0.118862$ & 1.572587 & 2.846171 & 1.646452\\
  11 & $-0.400698$ & 0.445410 & 0.030649 & 0.599906 & 2.303400 & 1.519685\\
  12 & 0.433222 & 0.413634 & 0.001792 & 0.598981 & 0.762273 & 1.567805\\
  13 & $-1.472310$ & $-0.460620$ & 0.044950 & 1.543336 & 3.444801 & 1.541667\\
  14 & 1.496445 & 0.941150 & 0.024802 & 1.767973 & 0.561416 & 1.556768\\
  15 & $-0.161600$ & 0.168879 & 0.012863 & 0.234094 & 2.334171 & 1.515821\\
  16 & $-0.697350$ & $-0.643470$ & $-0.128449$ & 0.957522 & 3.886829 & 1.705349
  \end{tabular}
  \end{center}

\item[(e)] Using radii calculated above and number of galaxies for each
  data one can calculate number of galaxies in each bin. Range of bin 1
  is (0.2, 0.425), bin 2 -- (0.425, 0.65) and so on bin 8 --
  (1.775, 2.0).

  Correct values for each bin are: (0.5p for each correct $n$)

  \begin{center}
  \begin{tabular}{rrr}
  & Center of the bin & Number of the galaxies\\
  1 & 0.312500 & 624\\
  2 & 0.537500 & 451\\
  3 & 0.762500 & 354\\
  4 & 0.987500 & 293\\
  5 & 1.212500 & 253\\
  6 & 1.437500 & 210\\
  7 & 1.662500 & 174\\
  8 & 1.887500 & 170
  \end{tabular}
  \end{center}
  \ioaafig{2022_DA_2-s4.pdf}

  In order to get probability density function each value of histogram
  must be divided by it integral on the whole region. The integral can
  be approximated by finding total area of the rectangles:
  \[ I = \sum_{n=1}^{8}H(r_i)\,d = 569.025 \]
  Where $d = 0.225$ is size of a bin. After dividing each data by $I$ we
  get (0.25p for each correct probability density):

  \begin{center}
  \begin{tabular}{rrrr}
  & Center of the bin & Number of the galaxies & Probability density\\
  1 & 0.312500 & 624 & 1.096613\\
  2 & 0.537500 & 451 & 0.792584\\
  3 & 0.762500 & 354 & 0.622117\\
  4 & 0.987500 & 293 & 0.514916\\
  5 & 1.212500 & 253 & 0.444620\\
  6 & 1.437500 & 210 & 0.369052\\
  7 & 1.662500 & 174 & 0.305786\\
  8 & 1.887500 & 170 & 0.298757
  \end{tabular}
  \end{center}
  \ioaafig{2022_DA_2-s5.pdf}

\item[(f)] Since histogram equation has following form:
  \[ H(r) = \frac{1}{A + B\cdot r} \]
  We should calculate $1/H(r)$ which gives the following values (0.125
  for each correct $1/H$):

  \begin{center}
  \begin{tabular}{rrrr}
  & $r$ & $H(r)$ & $1/H(r)$\\
  1 & 0.312500 & 624 & 0.001603\\
  2 & 0.537500 & 451 & 0.002217\\
  3 & 0.762500 & 354 & 0.002825\\
  4 & 0.987500 & 293 & 0.003413\\
  5 & 1.212500 & 253 & 0.003953\\
  6 & 1.437500 & 210 & 0.004762\\
  7 & 1.662500 & 174 & 0.005747\\
  8 & 1.887500 & 170 & 0.005882
  \end{tabular}
  \end{center}
  \ioaafig{2022_DA_2-s6.pdf}

  Finding the slope and the intercept of this line (by graph or any
  other method) gives following values:
  \[ A = 0.00065977;\qquad B = 0.00285494 \]
\end{description}

\textbf{Solutions -- Part 2}

\begin{description}
\item[(a)] We can $\phi$ values from Part 1.c to distribute them in 60
  degree bins.
  % sic: "We can phi values" as in source (missing "use")
  In bin 1 we have number of galaxies which fall in (0, 60) degree
  range, in bin 2 -- (60, 120) and so on bin 6 -- (300, 360).

  Part 1.c data in degrees:

  \begin{center}
  \begin{tabular}{rrr}
  & Angle & Number of galaxies\\
  1 & 302.496846 & 291\\
  2 & 353.752245 & 170\\
  3 & 303.884926 & 253\\
  4 & 290.179429 & 8\\
  5 & 115.183442 & 72\\
  6 & 294.241133 & 120\\
  7 & 133.773013 & 302\\
  8 & 73.435830 & 52\\
  9 & 203.897219 & 28\\
  10 & 163.073590 & 40\\
  11 & 131.975076 & 72\\
  12 & 43.675004 & 379\\
  13 & 197.372555 & 82\\
  14 & 32.166754 & 62\\
  15 & 133.738158 & 305\\
  16 & 222.698869 & 293
  \end{tabular}
  \end{center}

  We get the following data for the histogram (0.5p for each correct
  $n$):

  \begin{center}
  \begin{tabular}{rrr}
  & Center of the bin & Number of galaxies\\
  1 & 30 & 441\\
  2 & 90 & 124\\
  3 & 150 & 719\\
  4 & 210 & 403\\
  5 & 270 & 128\\
  6 & 330 & 714
  \end{tabular}
  \end{center}

  Uniform value is total number of galaxies divides by 6 which is
  421.5. Square deviation is $241^\circ$ (1p)
  \ioaafig{2022_DA_2-s7.pdf}

\item[(b)] A galaxy with the orbital radius $r$ around central mass $M$
  has velocity (2p):
  \[ v = \sqrt{\frac{GM}{r}} \]
  This will cause to give the following error in distance along line of
  sight (2p):
  \[ d = \frac{v}{H} = \frac{1}{H}\sqrt{\frac{GM}{r}} \]
  Angle difference for galaxies at 90 and 270 degrees will be (2p):
  \[ \alpha = \arctan\frac{d}{r}
     = \arctan\frac{1}{H}\sqrt{\frac{GM}{r^3}} \]
  \ioaafig{2022_DA_2-s8.pdf}

\item[(c)] Galaxies originally in bins of range (60, 120) and (240, 300)
  might move to (120, 180) and (300, 360) respectively. So we should
  subtract correction value found in previous task to the galaxies in
  bins 2 and 5 to see if they move to bins 1 and 4 respectively. (0.25p
  for each correct corrected angle)

  \begin{center}
  \begin{tabular}{rrrr}
  & Original angle & Corrected angle & Number of galaxies\\
  1 & 302.494665 & 276.136451 & 291\\
  2 & 353.752069 & 351.714594 & 170\\
  3 & 303.884728 & 300.199512 & 253\\
  4 & 290.178642 & 290.178642 & 8\\
  5 & 115.184577 & 115.184577 & 72\\
  6 & 294.241835 & 294.241835 & 120\\
  7 & 133.776280 & 125.061576 & 302\\
  8 & 73.437747 & 73.437747 & 52\\
  9 & 203.899065 & 203.899065 & 28\\
  10 & 163.073490 & 160.346665 & 40\\
  11 & 131.977481 & 120.548626 & 72\\
  12 & 43.676971 & 43.676971 & 379\\
  13 & 197.372956 & 197.372956 & 82\\
  14 & 32.167948 & 32.167948 & 62\\
  15 & 133.743596 & 94.077612 & 305\\
  16 & 222.694636 & 222.694636 & 293
  \end{tabular}
  \end{center}

  New corrected distribution looks like this (0.3p for each correct
  $n$):

  \begin{center}
  \begin{tabular}{rrr}
  & Center of the bin & Number of galaxies\\
  1 & 30 & 441\\
  2 & 90 & 429\\
  3 & 150 & 414\\
  4 & 210 & 403\\
  5 & 270 & 419\\
  6 & 330 & 423
  \end{tabular}
  \end{center}

  In this case square deviation equals to $12^\circ$ (1p)
  \ioaafig{2022_DA_2-s9.pdf}
\end{description}

\textbf{Solutions -- Part 3}

\begin{description}
\item[(a)] With simple combinatorics one can get (0.6p for each correct
  $N$):
  \begin{align*}
  N_{CC} &= N_C(N_C - 1) = 249500\\
  N^{*}_{CS} &= N_C N_S = 500000\\
  N_{CS} &= N_C(N_C - 1)N_S = 249500000\\
  N^{*}_{SS} &= N_C N_S(N_S - 1) = 499500000\\
  N_{SS} &= N_C N_S(N_C - 1)N_S = 249500000000
  \end{align*}

\item[(b)] We can calculate all probability densities using
  center-center and center-satellite densities. Non-normalized values
  (with $c_1 = c_2 = c_3 = 1$) are calculated in the following table
  (0.5p for each correct tuple):

  \begin{center}
  \begin{tabular}{rrrrrrr}
  & $r$ & $\rho^{*}_{CS}$ & $\rho_{CC}$ & $\rho^{*}_{SS}$
    & $\rho_{SS}$ & $\rho_{CS}$\\
  1 & 0.312500 & 1.096613 & 0.406141 & 0.375800 & 17.365137 & 0.975747\\
  2 & 0.537500 & 0.792584 & 0.807863 & 0.337652 & 12.992898 & 0.994072\\
  3 & 0.762500 & 0.622117 & 1.340510 & 0.295110 & 8.653487 & 0.988565\\
  4 & 0.987500 & 0.514916 & 1.134011 & 0.261824 & 8.568351 & 0.837480\\
  5 & 1.212500 & 0.444620 & 0.478630 & 0.239696 & 10.729180 & 0.638199\\
  6 & 1.437500 & 0.369052 & 0.143154 & 0.195787 & 10.112507 & 0.458355\\
  7 & 1.662500 & 0.305786 & 0.073165 & 0.155452 & 8.262279 & 0.352612\\
  8 & 1.887500 & 0.298757 & 0.060971 & 0.168470 & 8.998531 & 0.379962
  \end{tabular}
  \end{center}

  For normalization (same way as in Part 1.d) we get $c_1 = 2.18592$;
  $c_2 = 0.79304$; $c_3 = 0.0516$. Multiplying relevant densities by
  this values gives (0.5p for each correct normalized tuple):

  \begin{center}
  \begin{tabular}{rrrrr}
  & $r$ & $\rho^{*}_{SS}$ & $\rho_{SS}$ & $\rho_{CS}$\\
  1 & 0.312500 & 0.822854 & 0.900750 & 0.770962\\
  2 & 0.537500 & 0.739325 & 0.673957 & 0.785441\\
  3 & 0.762500 & 0.646175 & 0.448866 & 0.781089\\
  4 & 0.987500 & 0.573292 & 0.444450 & 0.661713\\
  5 & 1.212500 & 0.524839 & 0.556535 & 0.504257\\
  6 & 1.437500 & 0.428697 & 0.524548 & 0.362158\\
  7 & 1.662500 & 0.340380 & 0.428574 & 0.278607\\
  8 & 1.887500 & 0.368883 & 0.466764 & 0.300217
  \end{tabular}
  \end{center}
  \ioaafig{2022_DA_2-s10.pdf}
  \ioaafig{2022_DA_2-s11.pdf}
  \ioaafig{2022_DA_2-s12.pdf}
  \ioaafig{2022_DA_2-s13.pdf}

\item[(c)] Total probability density is weighted sum of all densities
  (4p):
  \[ \rho \sim \frac{N_{CC}\rho_{CC} + N^{*}_{CS}\rho^{*}_{CS}
     + N_{CS}\rho_{CS} + N^{*}_{SS}\rho^{*}_{SS} + N_{SS}\rho_{SS}}
     {N_{CC} + N^{*}_{CS} + N_{CS} + N^{*}_{SS} + N_{SS}} \]

  Calculated values and normalized densities are given in the following
  table:

  \begin{center}
  \begin{tabular}{rrrr}
  & $r$ & Calculated $\rho$ & Normalized $\rho$\\
  1 & 0.312500 & 0.004212 & 0.843662\\
  2 & 0.537500 & 0.003606 & 0.722433\\
  3 & 0.762500 & 0.002967 & 0.594357\\
  4 & 0.987500 & 0.002693 & 0.539524\\
  5 & 1.212500 & 0.002663 & 0.533361\\
  6 & 1.437500 & 0.002265 & 0.453652\\
  7 & 1.662500 & 0.001813 & 0.363239\\
  8 & 1.887500 & 0.001968 & 0.394216
  \end{tabular}
  \end{center}
  \ioaafig{2022_DA_2-s14.pdf}

  Most similar to this density is plot E on Figure 1 (1p).
\end{description}
